% %%
\documentclass[lettersize,journal]{IEEEtran}
%\documentclass[lettersize,journal,draft]{IEEEtran}
% This version has the human touch

\usepackage{amsmath,amsfonts,amssymb}
\usepackage{algorithmic}
\usepackage{array}
\usepackage{textcomp}
\usepackage{stfloats}
\usepackage{url}
\usepackage{verbatim}
\usepackage{graphicx}
\usepackage{cite}
\usepackage{balance}
\newtheorem{theorem}{Theorem}
\newtheorem{proposition}{Proposition}
\newtheorem{lemma}{Lemma}
\newtheorem{corollary}{Corollary}
\newtheorem{remark}{Remark}
\newtheorem{assumption}{Assumption}
\hyphenation{op-tical net-works semi-conduc-tor IEEE-Xplore}
\def\BibTeX{{\rm B\kern-.05em{\sc i\kern-.025em b}\kern-.08em
    T\kern-.1667em\lower.7ex\hbox{E}\kern-.125emX}}

\urlstyle{same}
\begin{document}

\title{Multirobot Tracking of
Instantaneous Coherent Structures}

\author{Christopher Waight and Christopher A. Kitts,~\IEEEmembership{Senior Member,~IEEE}
\thanks{Manuscript received Month Day, Year; revised Month Day, Year.}
\thanks{C. Waight is a PhD Candidate in the Robotic Systems Laboratory, Department of Mechanical
Engineering, Santa Clara University, Santa Clara, CA 95053 USA (e-mail: cwaight@scu.edu)}
\thanks{C. A. Kitts is Professor and Director of the Robotic Systems Laboratory, Department of
Mechanical Engineering, Santa Clara University, Santa Clara, CA 95053 USA (e-mail: ckitts@scu.edu)}
}
\markboth{Journal of \LaTeX\ Class Files,~Vol.~XX, No.~X, Month~Year}%
{Waight and Kitts: Multirobot Tracking of Instantaneous Coherent Structures}

\maketitle

% ---------------------------------------------------------------------------
\begin{abstract}
Short-term transport in a geophysical flow is organized by instantaneous
coherent structures: curves of strongest local attraction and repulsion,
along which drifting material collects within hours. Existing robotic
approaches to tracking such structures either commit to a presumed
structure before tracking begins or must integrate measurements over a
finite time horizon before an estimate exists. This paper presents a
multirobot method for finding and following these structures in real
time, with no map, no forecast, and no global information, using only
the flow velocity measured by each robot at its own position. A
distributed estimator recovers the local quadratic model of the field
from a single synchronous set of measurements, and six robots are
proved to be the minimum team for which this recovery is possible. The
same twelve fitted coefficients deliver two scalar surrogate fields
whose valley lines mark the structures: the Jacobian determinant $D$,
whose valleys trace the separatrices of the flow, and the smaller
rate-of-strain eigenvalue $s_1$, whose structures are invariant under
every time-dependent rotation of the observer. One adaptive navigation
controller per surrogate follows these valley lines. The $D$ tracker
rides the determinant's connected trench network on a modified Newton
step with a network-traversal guarantee; the $s_1$ tracker consumes
only frame-invariant quantities, and in a rotating-frame experiment it
recovers the same material curve in either observer frame while the $D$
tracker follows a frame artifact. Both controllers carry stability
certificates, and Monte Carlo noise sweeps place the $D$ tracker's
failure cliff exactly at the estimator's own reliability threshold, a
ceiling the $s_1$ tracker's first-order eigenframe defers by an order
of magnitude. Running unmodified on
ocean-radar currents from the Santa Barbara Channel, the $D$ tracker
follows a coherent ridge for 28 hours to a consistent landfall, and the
$s_1$ tracker converges to and holds a transient attracting profile
within one radar grid cell of ground truth computed offline from the
full map.
\end{abstract}

\begin{IEEEkeywords}
Multirobot Systems, Adaptive Navigation, Coherent Structures,
Okubo-Weiss, Rate of Strain, Objectivity, Cooperative
Estimation, Cluster Space Control
\end{IEEEkeywords}


% ---------------------------------------------------------------------------
\section{Introduction}
\label{sec:intro}

\IEEEPARstart{M}{ultirobot} systems exploit spatial distribution across a team to sample
and act on an environment more effectively than a single platform can,
with broad application to formation assembly, coverage, and distributed
sensing \cite{1}. Equipped with field sensors, such teams form mobile
sensor networks, and adaptive navigation (AN) techniques use their
real-time measurements to direct the network to or along features of
interest. Unlike survey-oriented navigation, which traverses a region
exhaustively and identifies features afterward, AN locates local
features quickly and follows them as they evolve.

AN in scalar fields is well developed. Each robot contributes a single
scalar measurement to a navigation policy that reaches and follows
features such as extrema, contour lines, ridges, trenches, saddle
points, and fronts \cite{26,27,28,29}, with formations sized so that
differential measurements across the cluster recover the gradient and
curvature information each policy consumes \cite{15,26}.

Some environments are better modeled by vector fields: ocean currents,
atmospheric jets, and other geophysical flows assign a velocity, not a
scalar, to each point. The organizing features of such fields begin
with their critical points, and in our own prior work a three-robot
cluster estimated the location and type of an isolated critical point
from a linear fit of simultaneous local velocity measurements
\cite{14}, a point-estimation problem whose target is a single
coordinate. Coordinated multi-vehicle sampling of dynamic ocean
features is itself an active area of marine robotics, from
drifter-referenced AUV surveys of advecting patches \cite{2} to
heterogeneous ASV/AUV teams tracking salinity fronts \cite{3} and
glider fleets adaptively sampling upwelling events during the AOSN-II
experiment in Monterey Bay \cite{4}.

Beyond isolated critical points, these missions increasingly call for
locating and tracking dynamically meaningful boundaries. The boundary
that appears prominently in steady two-dimensional flows is the
separatrix, which contains the stable and unstable manifolds connecting
to saddle points and partitions the domain into basins whose
trajectories do not mix. Tracking it is relevant to dispersion
modeling, pollutant containment, and the deployment of sensor networks
in coastal and atmospheric environments \cite{6,7,8}, and, like the
related Okubo-Weiss field used to detect and classify oceanic mesoscale
eddies \cite{5}, it can in principle be inferred from instantaneous
field measurements alone.

A line of work has equipped robot teams to track a coherent structure
directly in the flow rather than reconstruct it offline, by straddling a
presumed manifold with a small formation \cite{9,10,11,12} or by
accumulating an onboard FTLE estimate over a finite time horizon
\cite{13}. As Section~\ref{sec:related_robots} details, each of these
either commits to the structure's identity before tracking begins or
pays an integration latency before an estimate exists.

This paper's approach requires neither a seeded manifold nor an
integration horizon. From an arbitrary six-robot configuration, the
team infers a structure's location, its type, and, at a crossing, which
branch to continue onto, from the current instantaneous measurement
alone using a second-order estimate. The technique exploits the fact
that the velocity gradient of a flow splits uniquely into strain plus
spin, $\mathbf{J} = \mathbf{S} + \mathbf{W}$, and that for
incompressible planar flow the determinant splits with it, $D =
\det(\mathbf{J}) = \omega^2/4 - s_1^2$, where $\omega$ is the vorticity
carried by $\mathbf{W}$ and $s_1$ the smaller eigenvalue of $\mathbf{S}$
(Section~\ref{sec:detj_field}). Each term of the splitting is a
trackable surrogate of the flow's transport skeleton, and each of the
two controllers developed here rides one of them.

This paper presents the following contributions:
\begin{enumerate}
    \item A distributed second-order estimator: six robots are proved
          to be the exact minimum for recovering the quadratic model of
          a coupled two-component field, the same team size with which
          \cite{15} recovers the six coefficients of a scalar field,
          here recovering twelve, and one synchronous sample yields
          closed forms for $D$, $\nabla D$, $\mathbf{H}_D$ and for
          $s_1$, $\nabla s_1$, and the strain eigenframe, with exact,
          heading-isotropic noise gains.
    \item The $D$ tracker: a hybrid controller built on a
          modified Newton step whose eigen-directions are independently
          saturated and sign-flipped, so the team snaps onto a trench
          of $D$, slides along it, and crosses crests under flow
          assistance. A network-traversal theorem guarantees passage
          through the saddle points where trench segments cross,
          committing to no manifold identity in advance (unlike
          \cite{9,10}), with terminal capture as a tightened-threshold
          special case.
    \item The $s_1$ tracker: an objective controller built from the
          strain eigenframe, whose design is dictated by a structural
          impossibility result (the $s_1$ Hessian of any quadratic
          velocity fit is negative semidefinite), together with a
          closed-loop frame-equivariance proposition.
    \item Validation of both trackers under a fully specified noise
          model: Monte Carlo sweeps placing the $D$ tracker's failure
          cliff at the estimator's own reliability threshold and the
          $s_1$ tracker's an order of magnitude beyond it, a rotating-frame
          experiment demonstrating objectivity in closed loop, and runs
          on real Santa Barbara Channel radar currents. Hardware
          validation is future work (Section~\ref{sec:future}).
\end{enumerate}

The remainder of the paper is organized as follows. The rest of this
section reviews related work. Section~\ref{sec:detj_field} formulates
the problem: the two surrogate fields, the splitting identity that
relates them, and the double-gyre running example.
Section~\ref{sec:estimation} develops the six-robot quadratic estimator
and its accuracy under noise. Section~\ref{sec:controllers} presents
the two adaptive navigation controllers and their stability properties,
with proofs in Appendix~\ref{app:proofs}. Section~\ref{sec:results}
describes the simulation program and presents results, including the
ocean experiments. Section~\ref{sec:discussion} discusses implications,
failure modes, and limitations. Section~\ref{sec:conclusion} concludes.


% ---------------------------------------------------------------------------

\subsection{Related Work}
\label{sec:related_robots}

\subsubsection{Coherent Structures and Instantaneous Criteria}
Haller and Yuan formalized the Lagrangian coherent structure (LCS) as
the generalization of stable and unstable manifolds to flows with
arbitrary time dependence \cite{6}, with the standard computable
definition, a ridge of the finite-time Lyapunov exponent (FTLE) field,
given by Shadden, Lekien, and Marsden on the same double-gyre model
used throughout this paper \cite{8}. Haller's review makes objectivity,
invariance under time-dependent Euclidean frame changes, the admission
test for physical meaning, and excludes the classical instantaneous
criteria, Okubo-Weiss among them, for lacking it \cite{16}. Lagrangian
definitions pay for objectivity with latency: an FTLE ridge requires
the field's future evolution, integrated over a finite horizon. Serra
and Haller closed the gap from the instantaneous side with objective
Eulerian coherent structures (OECSs), built from the rate-of-strain
tensor and shown to be the zero-horizon limits of LCSs \cite{17}, a
limit made precise for the FTLE field in \cite{18}; velocity-gradient
recovery from sparse trajectories \cite{19} and the strain-spin
decomposition itself \cite{20} remain active. The instantaneous
criteria nevertheless remain the computable workhorses: Okubo and Weiss
partition a planar flow from the local velocity gradient alone
\cite{21,22}, generalized to three dimensions \cite{23}, sharpened
oceanographically \cite{24}, and used at scale in mesoscale eddy census
work \cite{5}. Unlike approaches that sense a scalar reduction of the
field directly, which cannot tell a sink from a saddle from a vortex
\cite{25}, the team here fits the full local quadratic vector field and
obtains $D$ and $s_1$, with their gradients and eigenstructure, as
derived quantities: the surrogate's shape, not merely its sign, is what
the controllers consume.

\subsubsection{Robotic Tracking of Flow Structures}
The nearest prior work tracks a coherent structure directly with a
robot formation. Michini et al. bracket a presumed stable or unstable
manifold with three robots using only local velocity measurements
\cite{9}; the strategy extends to $N$ robots \cite{10}, was validated
on micro autonomous surface vehicles in a flow tank \cite{11}, and
related flow-based control operates marine robots in gyre-like
environments \cite{12}. The straddle family commits to the manifold's
identity before tracking begins: the branch is chosen in advance and
never revisited, the validity proof assumes the segment crosses the
manifold exactly once, and neither controller nor proof addresses the
saddle point itself, exactly where that identity becomes ambiguous.
Kularatne and Hsieh instead compute the local FTLE field on board,
recovering the structure's type at the price of a nonzero integration
horizon before an estimate exists \cite{13}. None of these strategies
estimates the curvature of the field itself, which is what fixes not
just where a structure lies but what kind it is and how many branches
meet there; the estimator of Section~\ref{sec:estimation} supplies
exactly that, so the branch to continue onto at a crossing is measured
fresh each control cycle rather than assumed.

\subsubsection{Cooperative Field Estimation}
On the estimation side, \"{O}gren, Fiorelli, and Leonard established
cooperative gradient climbing with a formation of mobile sensors
\cite{26}, and Bri\~{n}\'{o}n-Arranz, Renzaglia, and Schenato extended
the idea to the Hessian: five robots on a ring plus one at the center
recover the gradient and Hessian of a scalar signal in closed form
\cite{15}. McDonald, Kitts, and Neumann developed the complementary
adaptive-navigation primitives that follow scalar-field ridges,
trenches, and saddle points from local differences alone
\cite{27,28,29}. Every prior formation samples one signal channel;
recovering the Jacobian and its curvature for a coupled two-component
field doubles the estimation problem to twelve unknowns, met by the two
measurements each robot contributes, and Section~\ref{sec:estimation}
proves six robots, the same count as the ring-plus-center of \cite{15},
is the exact minimum for the coupled case, for any non-conic placement
and against any algorithm. Deriving $D$ and $s_1$ by differentiating through the
jointly fitted vector coefficients, rather than fitting either
surrogate as an independent scalar surface, preserves the coupling
between $u$ and $v$ that a scalar-surrogate route would discard.

\subsubsection{Cluster Space Control and Testbeds}
The team itself is controlled in cluster space: Kitts and Mas
parameterize a formation by aggregate quantities and invert that
parameterization analytically to distribute one formation-level command
to every robot \cite{30}, closed in the loop for gradient-based
navigation on physical surface vessels in \cite{31}. The
side-angle-side (SAS) parameterization used for vector-field navigation
in \cite{14} is one instance, and the six-robot formation here is built
from three SAS pairs sharing a centroid
(Appendix~\ref{app:kinematics}); the grouping is invisible to the
estimator. The physical testbed behind this lineage evolved from a
low-cost twelve-robot indoor platform \cite{32} to the Decabot fleet
with OptiTrack motion capture \cite{33}, whose identified robot
dynamics this paper's simulations adopt; finite-time coordination
results such as \cite{34} pursue convergence-rate analyses that this
paper trades for geometric traversal guarantees. Where \cite{14} used a
linear fit with no curvature and \cite{15} used the raw Hessian inverse
to reach a scalar extremum, the modified Newton step introduced here
combines per-eigendirection saturation with the saddle-free sign flip
of Dauphin et al. \cite{35}, previously unapplied to a robot team
tracking a spatial field.


% ---------------------------------------------------------------------------
\section{Problem Formulation}
\label{sec:detj_field}

\subsection{Flow Fields and Local Polynomial Models}

A planar flow assigns a velocity vector $\mathbf{v}(\mathbf{p}) =
(u(\mathbf{p}), v(\mathbf{p}))$ to each position $\mathbf{p} = (x, y)$ in its
domain. Its Jacobian at $\mathbf{p}$ is
\begin{equation}
    \mathbf{J}(\mathbf{p}) =
    \begin{bmatrix}
        \partial u / \partial x & \partial u / \partial y \\
        \partial v / \partial x & \partial v / \partial y
    \end{bmatrix},
    \label{eq:jacobian}
\end{equation}
and we write $D(\mathbf{p}) = \det(\mathbf{J}(\mathbf{p}))$ throughout.

The surrogate fields this paper tracks are built from $\mathbf{J}$, and
navigating them requires their first and second spatial derivatives,
which are polynomial functions of the first and second derivatives of
$u$ and $v$. The team therefore fits a local quadratic model of each
component,
\begin{align}
    u(\mathbf{p}_c + \tilde{\mathbf{p}}) &\approx u_0 + u_x \tilde{x}
        + u_y \tilde{y} + \tfrac{1}{2} u_{xx} \tilde{x}^2
        + u_{xy} \tilde{x}\tilde{y} + \tfrac{1}{2} u_{yy} \tilde{y}^2,
        \label{eq:quad_model_u} \\
    v(\mathbf{p}_c + \tilde{\mathbf{p}}) &\approx v_0 + v_x \tilde{x}
        + v_y \tilde{y} + \tfrac{1}{2} v_{xx} \tilde{x}^2
        + v_{xy} \tilde{x}\tilde{y} + \tfrac{1}{2} v_{yy} \tilde{y}^2,
        \label{eq:quad_model_v}
\end{align}
where $\tilde{\mathbf{p}} = (\tilde{x}, \tilde{y})$ is position relative to
the cluster centroid $\mathbf{p}_c$ and the twelve coefficients are the
Taylor data of the field at $\mathbf{p}_c$: the value, gradient, and
Hessian of each component. 

\subsection{Separatrices and Eulerian Surrogates}
\label{sec:separatrix}

At a saddle point of a steady flow, the stable and unstable manifolds
are the curves of initial conditions that converge to the saddle in
forward and backward time; together they form the separatrix, which
partitions the domain into regions whose trajectories do not mix. In
time-dependent flows the analogous objects are LCSs, ridges of the FTLE
field \cite{6,7,8}, but computing FTLE requires integrating dense grids
of trajectories over a finite horizon, which presumes exactly the
global, forward-in-time knowledge this paper's sensing model excludes.
The alternative is an Eulerian surrogate: a quantity computed pointwise
from the instantaneous field and its spatial derivatives. This paper
tracks two, both delivered by the same fit. The determinant field $D$
is Galilean invariant, globally connected, and coincides exactly with
the separatrix for the steady double gyre
(Section~\ref{sec:dg_example}); it is evaluated against FTLE fields
computed offline from ocean radar data in Section~\ref{sec:results}.
The strain eigenvalue $s_1$ underlies the attracting OECSs of
\cite{17,18} and is fully objective. The two are the terms of the
splitting identity developed in Section~\ref{sec:strain_field}, and the
price and payoff of each choice are quantified in
Section~\ref{sec:results}.

\subsection{Surrogate I: the Okubo-Weiss Field $D$}

Write the entries of $\mathbf{J}$ as in (\ref{eq:jacobian}). The normal
strain rate, shear strain rate, and vorticity are
$s_n = \partial u/\partial x - \partial v/\partial y$,
$s_s = \partial v/\partial x + \partial u/\partial y$, and
$\omega = \partial v/\partial x - \partial u/\partial y$. The Okubo-Weiss
parameter \cite{21,22} is $Q = \tfrac{1}{4}(s_n^2 + s_s^2 - \omega^2)$, and
direct expansion gives the algebraic identity
\begin{equation}
    Q = \tfrac{1}{4}\bigl(\nabla \cdot \mathbf{v}\bigr)^2 - D.
    \label{eq:okubo_weiss}
\end{equation}
For the incompressible flows considered in this paper,
$\nabla \cdot \mathbf{v} = 0$ and $Q = -D$ exactly, so the Okubo-Weiss
field and the Jacobian determinant carry the same information with
opposite sign. We work with $D$.

The sign of $D$ fixes the local eigenstructure. Incompressibility gives
$\operatorname{tr}(\mathbf{J}) = 0$, so the eigenvalues of $\mathbf{J}$
satisfy $\lambda^2 = -D$:
\begin{equation}
    \lambda = \pm\sqrt{-D} \ \ (D < 0), \qquad
    \lambda = \pm i\sqrt{D} \ \ (D > 0).
    \label{eq:eig_detj}
\end{equation}
Where $D < 0$ the flow is locally hyperbolic (strain-dominated) and
nearby particles separate at the instantaneous exponential rate
$\sqrt{-D}$; where $D > 0$ it is locally elliptic (rotation-dominated).
The convention used throughout this paper is therefore: $D > 0$ marks
vortex cores, $D < 0$ marks strain regions, and the zero level set
\begin{equation}
    \mathcal{C} = \{\mathbf{p} : D(\mathbf{p}) = 0\}
    \label{eq:ow_boundary}
\end{equation}
is the Okubo-Weiss boundary between the two regimes. By
(\ref{eq:eig_detj}), the more negative $D$ is, the stronger the local
hyperbolicity, so valley lines (trenches) of $D$ inside strain regions
mark the loci of strongest instantaneous attraction and repulsion. The
boundary $\mathcal{C}$ delimits the strain region within which the
trench, the Eulerian structure this paper tracks, is found.

\subsection{Surrogate II: the Strain Eigenvalue $s_1$}
\label{sec:strain_field}

The velocity gradient decomposes into symmetric and antisymmetric
parts, $\mathbf{J} = \mathbf{S} + \mathbf{W}$, with rate-of-strain
tensor $\mathbf{S} = (\mathbf{J} + \mathbf{J}^\top)/2$ and spin tensor
$\mathbf{W} = (\mathbf{J} - \mathbf{J}^\top)/2$, which carries the
vorticity. Under a time-dependent Euclidean change of observer
$\mathbf{p}' = \mathbf{Q}(t)\mathbf{p} + \mathbf{b}(t)$, the velocity
gradient transforms as $\mathbf{J}' = \mathbf{Q}\mathbf{J}
\mathbf{Q}^\top + \dot{\mathbf{Q}}\mathbf{Q}^\top$. The frame's own
rotation $\dot{\mathbf{Q}}\mathbf{Q}^\top$ is antisymmetric, so it
lands entirely in the spin part, $\omega' = \omega + 2\Omega$ for a
frame spinning at rate $\Omega$, while the strain part transforms
cleanly, $\mathbf{S}' = \mathbf{Q}\mathbf{S}\mathbf{Q}^\top$. The
eigenvalues $s_1 \leq s_2$ of $\mathbf{S}$ are therefore
frame-invariant, its eigenvectors $\mathbf{e}_1, \mathbf{e}_2$
co-rotate with the frame, and quantities built from them are objective:
material behavior, which cannot depend on the observer, can be
predicted from them in any frame \cite{16,17}.

The two surrogates of this paper are the two halves of this splitting.
For any symmetric $\mathbf{S}$ and antisymmetric $\mathbf{W}$ the cross
term vanishes, $\operatorname{tr}(\mathbf{S}\mathbf{W}) = 0$, so the
determinant splits exactly:
\begin{equation}
    D \;=\; \det\mathbf{S} + \det\mathbf{W}
      \;=\; s_1 s_2 + \frac{\omega^2}{4}
      \;=\; \frac{\omega^2}{4} - s_1^2,
    \label{eq:splitting}
\end{equation}
the last equality by incompressibility ($s_2 = -s_1$). Three
consequences of (\ref{eq:splitting}) organize everything that follows.
First, the Okubo-Weiss partition is the competition between the two
terms: rotation-dominated where the spin term wins ($D > 0$),
strain-dominated where the strain term wins ($D < 0$). Second,
substituting $\omega' = \omega + 2\Omega$ gives
\begin{equation}
    D' = D + \Omega\,\omega + \Omega^2,
    \label{eq:D_not_objective}
\end{equation}
so two observers rotating relative to one another disagree about where
the $D$ structures are: the frame artifact lives entirely in the spin
term, and only there. Third, where the vorticity vanishes the identity
collapses to $s_1 = -\sqrt{-D}$: the two surrogates coincide exactly on
vorticity-free lines and part ways everywhere else, and in every
rotating frame.

The smaller eigenvalue $s_1$ is the local instantaneous contraction
rate: where $s_1 < 0$, material filaments are compressed along
$\mathbf{e}_1$ and stretched along $\mathbf{e}_2$ at rate $s_2 = -s_1$.
Serra and Haller's attracting OECSs, operationalized as transient
attracting profiles (TRAPs) for search-and-rescue prediction
\cite{17,36}, are built from exactly this structure: a TRAP core is a
local minimum of $s_1$ with $s_1 < 0$, and the attracting curve through
the core is tangent to the stretching eigenvector $\mathbf{e}_2$ with
normal $\mathbf{e}_1$, terminated where attraction stops dominating. In
the language already introduced, attracting OECSs are trenches of the
signed scalar field $s_1$, ridden along $\mathbf{e}_2$; the objective
tracker of Section~\ref{sec:oecs_controller} does to $s_1$ what the
$D$ tracker does to $D$, with one structural difference in what
the estimator can supply (Section~\ref{sec:s1_estimation}). One caution
accompanies the eigenstructure: $\mathbf{e}_1, \mathbf{e}_2$ form
direction fields with no globally consistent sign, undefined at
isotropic points of $\mathbf{S}$ ($s_1 = s_2$), where tensor lines
exhibit half-integer-index singularities; the controller must therefore
orient its tangent by continuity rather than by any pointwise rule.

\subsection{Running Example: the Double Gyre}
\label{sec:dg_example}

The double-gyre flow used throughout is the steady ($\varepsilon = 0$)
member of the Shadden family \cite{8}, shifted from the canonical domain
$[0,2] \times [0,1]$ to $[-1,1] \times [-0.5, 0.5]$ by
$x_f = x + 1$, $y_f = y + 0.5$, where $(x_f, y_f)$ are field-frame and
$(x, y)$ world-frame coordinates. The field is
\begin{equation}
    u(\mathbf{p}) = -\pi A \sin(\pi x_f) \cos(\pi y_f), \label{eq:dg_u}
\end{equation}
\begin{equation}
    v(\mathbf{p}) = \pi A \cos(\pi x_f) \sin(\pi y_f), \label{eq:dg_v}
\end{equation}
with $A = 0.1$. The flow is incompressible, the separatrix is the line
$x = 0$, the saddle points are $\mathbf{p}^*_1 = (0, -0.5)$ and
$\mathbf{p}^*_2 = (0, 0.5)$, and the gyre centers are $(\pm 0.5, 0)$.

All quantities of interest have closed forms (derivations in
Appendix~\ref{app:analytic}):
\begin{equation}
    D(\mathbf{p}) = -\frac{\pi^4 A^2}{2}
    \bigl(\cos 2\pi x_f + \cos 2\pi y_f\bigr),
    \label{eq:det_analytic_main}
\end{equation}
\begin{equation}
    \begin{aligned}
    \nabla D &= \pi^5 A^2
    \begin{bmatrix} \sin 2\pi x_f \\ \sin 2\pi y_f \end{bmatrix},
    \\
    \mathbf{H}_D &= 2\pi^6 A^2
    \begin{bmatrix} \cos 2\pi x_f & 0 \\ 0 & \cos 2\pi y_f \end{bmatrix}.
    \end{aligned}
    \label{eq:gradhess_analytic_main}
\end{equation}
Three facts follow directly and organize everything that comes after.

First, the Okubo-Weiss boundary is a diamond: since $\cos 2\pi x_f +
\cos 2\pi y_f = 2\cos\pi(x_f + y_f)\cos\pi(x_f - y_f)$, the zero set
$\mathcal{C}$ consists of the straight lines $x_f \pm y_f \in
\tfrac{1}{2} + \mathbb{Z}$, a diamond inscribed around each gyre center
with $D > 0$ inside.

Second, the separatrix is a trench of $D$, and the trench is one line
of a periodic grid. On $x = 0$ (that is, $x_f = 1$) the cross-track
derivative $\partial D/\partial x$ vanishes identically and the
cross-track curvature $2\pi^6 A^2$ is positive, so the separatrix is an
exact valley line of the signed field $D$
(Fig.~\ref{fig:trench_cross}); the vorticity also vanishes there, so by
(\ref{eq:splitting}) $\sqrt{-D}$ is exactly the instantaneous
attraction rate toward the separatrix. Restricted to the separatrix,
$D(0, y) = -\tfrac{\pi^4 A^2}{2}(1 + \cos 2\pi y_f)$ is deepest at the
two saddles and rises to a crest ($D = 0$) at the origin, where the two
diamonds touch: a controller that only descends $D$ along the trench
stalls at the crest, which motivates the flow-assisted mode of
Section~\ref{sec:sep_controller}. (The ambient flow on the separatrix
is tangent to it, from the upper saddle toward the lower, vanishing
only at the saddles.) Nor does the trench end at a saddle: by
periodicity every line $x_f \in \mathbb{Z}$ and every line $y_f \in
\mathbb{Z}$ is a valley of $D$ in its own cross-track coordinate, the
valley lines form an infinite grid, and the saddle points are exactly
its right-angle crossings. A tracker that reaches a saddle and
continues is still on a trench, now running perpendicular to the one it
arrived on; Section~\ref{sec:stab_sep} states the traversal result, and
Section~\ref{sec:results} reports which behavior, capture or
continuation, the estimator actually produces.

Third, the strain field shares the grid but signs its segments. The
double gyre's shear strain vanishes identically, so $\mathbf{S}$ is
diagonal everywhere and (Appendix~\ref{app:analytic})
\begin{equation}
    s_1 = -\pi^2 A\,\lvert\cos \pi x_f \cos \pi y_f\rvert,
    \label{eq:s1_analytic_main}
\end{equation}
with eigenvectors fixed on the coordinate axes away from degeneracies.
The trenches of $s_1$ form the same grid as the trenches of $D$, and
the TRAP cores, the minima $s_1 = -\pi^2 A$, sit at the six flow
saddles on the walls. But the attracting/repelling identity of the
segments alternates: the upper half of the separatrix is an attracting
OECS, the lower half is repelling (its attracting partner is the bottom
wall), and segments meet at the degeneracy lines $x = \pm 0.5$ and $y =
0$, where $\mathbf{S} = \mathbf{0}$. The crest of the $D$ trench at the
origin is thus also an isotropic point of $\mathbf{S}$, where the
attracting segment terminates in the Serra-Haller sense: an objective
tracker riding the upper separatrix should not sail through the crest,
because the structure it is tracking ends there. The vanishing shear
also makes the double gyre nongeneric, with degeneracy lines in place
of the isolated singular points of a generic strain field \cite{17},
which the ocean field of Section~\ref{sec:ocean_sim} does exhibit.

\begin{figure}[!t]
    \centering
    \includegraphics[width=0.45\textwidth]{figures/detJ_trench_cross_section.png}
    \caption{Signed $D = \det(\mathbf{J})$ as an elevation surface over the
    double-gyre domain. The two gyre cores rise as ridges ($D > 0$,
    rotation-dominated) and the exterior falls below zero ($D < 0$,
    strain-dominated); the dashed curve on the floor is the $D = 0$ boundary.
    The separatrix at $x = 0$ (bold) is a valley of $D$ running between the two
    ridges, its floor rising to $D = 0$ at the origin and deepening toward the
    saddle points.}
    \label{fig:trench_cross}
\end{figure}

\subsection{Problem Statement}
\label{sec:problem_statement}

Consider $N$ robots at positions $\mathbf{p}_i$, each measuring the local
flow $(u_i, v_i) = \mathbf{v}(\mathbf{p}_i)$ at the common sampling
instant. No map, forecast, or prior model of the field is available, and
no robot can integrate trajectories of the field. Using only the current
synchronous samples $\{(\mathbf{p}_i, u_i, v_i)\}_{i=1}^{N}$ at each
control cycle, the team must generate a centroid velocity command
$\dot{\mathbf{p}}_c$ that solves one of two problems, one per term of
(\ref{eq:splitting}):

\emph{Problem 1 (separatrix network traversal).} Drive the cluster
centroid onto the trench of $D$ and travel it, in the direction of the
ambient flow, through the saddle points where trench segments cross;
terminal capture at a single saddle is the special case in which the
selector's stationary-point test is tightened to stop the team there.

\emph{Problem 2 (objective structure tracking).} Drive the centroid
onto the attracting OECS nearest its start using only objective
quantities, and either park at its TRAP core or ride it until the
structure ends.

Both require, at minimum, pointwise estimates of the surrogate, its
gradient, and its curvature or eigenframe at the centroid. The next
section shows how a six-robot formation produces all of them from one
fit.


% ---------------------------------------------------------------------------
\section{Distributed Second-Order Estimation}
\label{sec:estimation}

\subsection{Quadratic Fit from Six Simultaneous Measurements}

Let $\tilde{\mathbf{p}}_i = \mathbf{p}_i - \mathbf{p}_c$ denote robot
positions relative to the centroid and define the quadratic basis vector
\begin{equation}
    \boldsymbol{\phi}(\tilde{\mathbf{p}}) =
    \bigl[\,1, \; \tilde{x}, \; \tilde{y}, \; \tfrac{1}{2}\tilde{x}^2, \;
    \tilde{x}\tilde{y}, \; \tfrac{1}{2}\tilde{y}^2\,\bigr]^\top,
    \label{eq:basis}
\end{equation}
whose one-half factors make the fitted coefficients equal to the Taylor
derivatives in (\ref{eq:quad_model_u})--(\ref{eq:quad_model_v}). Stacking
the basis vectors of the six robots row-wise gives the $6 \times 6$
formation matrix $\boldsymbol{\Phi}$, the second-order generalization of
the $3 \times 3$ matrix $\mathbf{A}$ of \cite{14}. The coefficient vectors
$\boldsymbol{\theta}_u = [u_0, u_x, u_y, u_{xx}, u_{xy}, u_{yy}]^\top$ and
$\boldsymbol{\theta}_v = [v_0, v_x, v_y, v_{xx}, v_{xy}, v_{yy}]^\top$
solve the pair of linear systems
\begin{equation}
    \boldsymbol{\Phi}\,\boldsymbol{\theta}_u = \mathbf{u}, \qquad
    \boldsymbol{\Phi}\,\boldsymbol{\theta}_v = \mathbf{v},
    \label{eq:solve}
\end{equation}
where $\mathbf{u} = [u_1, \ldots, u_6]^\top$ and $\mathbf{v} = [v_1,
\ldots, v_6]^\top$ collect the measurements. One matrix factorization
serves both components. When the field is exactly quadratic over the
formation footprint, the solution of (\ref{eq:solve}) is exact for any
nonsingular $\boldsymbol{\Phi}$; for general smooth fields the model
error is quantified in Section~\ref{sec:sensitivity}.

Six robots are exactly the minimum team for this recovery. Each robot
contributes one linear equation per component to (\ref{eq:solve}), and
each system has six unknowns, so no algorithm, linear or otherwise, can
uniquely recover the quadratic model from one synchronous sample of
fewer than six robots (Proposition~\ref{prop:minimality},
Appendix~\ref{app:proofs}). Six suffice exactly when
$\boldsymbol{\Phi}$ is nonsingular, which fails precisely when the six
sampling points lie on a common conic, including degenerate conics such
as line pairs (Proposition~\ref{prop:conic}); five robots on a ring
with a sixth at the center never do (Corollary~\ref{cor:pentagon}). The
conic criterion carries a practical warning: any six robots on one
circle, a regular hexagon for example, are singular regardless of size
or orientation, a failure mode easy to adopt by accident because rings
are natural formations for mobile robots; moving one robot to the
center removes it entirely.

The pentagon-plus-center shape (Fig.~\ref{fig:pentagon_formation}) is
the same minimal configuration that Bri\~{n}\'{o}n-Arranz et al.\
derived for Hessian estimation of a scalar field \cite{15}. Their
estimator combines the ring measurements with fixed symmetric weights,
algebraically equivalent to solving (\ref{eq:solve}) on the ideal
formation, but the fixed weights assume that geometry and acquire a
bias proportional to the deformation the formation always exhibits
under real formation control. Solving (\ref{eq:solve}) at the measured
positions instead remains exact for quadratic fields under any
deformation short of the conic degeneracy of
Proposition~\ref{prop:conic}, so formation error affects accuracy only
through the conditioning of $\boldsymbol{\Phi}$, monitored online
through $\det(\boldsymbol{\Phi})$.

\begin{figure}[!t]
    \centering
    \includegraphics[width=0.30\textwidth]{figures/pentagon_formation.png}
    \caption{Six-robot formation: five robots on a ring of radius $\rho$
    at $72^\circ$ spacing and one robot at the center, the minimal
    symmetric formation for second-order estimation
    (Corollary~\ref{cor:pentagon}, \cite{15}). For formation control the six
    robots are grouped into three pairs whose midpoints define an SAS
    triangle (Appendix~\ref{app:kinematics}); the grouping is invisible to
    the estimator, which uses only the six measured positions.}
    \label{fig:pentagon_formation}
\end{figure}

\subsection{Noise Model}
\label{sec:noise_model}

Two independent noise channels perturb the samples available to the
estimator. Measurement noise is additive Gaussian on the $(u, v)$
components read by each robot:
\begin{equation}
    u_i^m = u(\mathbf{p}_i) + \eta_i^u, \quad
    v_i^m = v(\mathbf{p}_i) + \eta_i^v,
    \quad \eta \sim \mathcal{N}(0, \sigma_{uv}^2),
    \label{eq:meas_noise}
\end{equation}
with samples independent across robots, components, and time steps; this
is the sensor model of \cite{9}. Position noise perturbs the location at
which each robot samples the field,
\begin{equation}
    (u_i^m, v_i^m) = \mathbf{v}(\mathbf{p}_i + \boldsymbol{\xi}_i), \quad
    \boldsymbol{\xi}_i \sim \mathcal{N}(\mathbf{0}, \sigma_p^2 \mathbf{I}),
    \label{eq:pos_noise}
\end{equation}
while the quadratic fit (\ref{eq:solve}) uses the nominal positions
$\mathbf{p}_i$: the robot reads the field at a slightly wrong place. To
first order this induces an effective measurement error
$\mathbf{J}(\mathbf{p}_i)\,\boldsymbol{\xi}_i$, which couples the two
channels through the local Jacobian magnitude and is tested directly by
the two-dimensional noise sweep of Section~\ref{sec:test_plan}.
% Implementation: _sample_vector_at_robots in pentagon_primitives.py reads
% cluster.measurement_noise_std (eq:meas_noise) and
% cluster.position_noise_std (eq:pos_noise).

\subsection{Closed-Form Recovery of $D$, $\nabla D$, and $\mathbf{H}_D$}

The estimated Jacobian at the centroid is assembled from the linear
coefficients,
\begin{equation}
    \hat{\mathbf{J}} =
    \begin{bmatrix} u_x & u_y \\ v_x & v_y \end{bmatrix},
    \qquad
    \hat{D} = u_x v_y - u_y v_x.
    \label{eq:J_hat}
\end{equation}
Under the quadratic model, each entry of $\mathbf{J}$ is an affine
function of position, so $D(\mathbf{p})$ is itself an exactly quadratic
polynomial over the footprint. Its gradient and Hessian at the centroid
are polynomial functions of the fitted coefficients:
\begin{equation}
    \hat{\mathbf{g}} = \nabla \hat{D} =
    \begin{bmatrix}
        u_{xx} v_y + u_x v_{xy} - u_{xy} v_x - u_y v_{xx} \\
        u_{xy} v_y + u_x v_{yy} - u_{yy} v_x - u_y v_{xy}
    \end{bmatrix},
    \label{eq:grad_det}
\end{equation}
\begin{equation}
    \hat{\mathbf{H}}_D =
    \begin{bmatrix}
        2(u_{xx} v_{xy} - u_{xy} v_{xx}) &
        u_{xx} v_{yy} - u_{yy} v_{xx} \\
        u_{xx} v_{yy} - u_{yy} v_{xx} &
        2(u_{xy} v_{yy} - u_{yy} v_{xy})
    \end{bmatrix}.
    \label{eq:hess_det}
\end{equation}
These follow by differentiating $D = u_x v_y - u_y v_x$ with the
third-order terms absent, and require no field evaluations beyond the six
measurements. The omission is harmless for the gradient but not for the
Hessian; its systematic consequences are quantified in
Section~\ref{sec:sensitivity}. The estimated flow at the centroid,
$\mathbf{f} = (u_0, v_0)$, equals the center robot's reading
(Lemma~\ref{lem:gains}, Section~\ref{sec:sensitivity}). Together
(\ref{eq:J_hat})--(\ref{eq:hess_det}) supply every quantity the
controller of Section~\ref{sec:controllers} consumes: $\hat{D}$,
$\hat{\mathbf{g}}$, $\hat{\mathbf{H}}_D$, and $\mathbf{f}$.

\subsection{Closed-Form Recovery of $s_1$ and the Strain Eigenframe}
\label{sec:s1_estimation}

The same fit supplies the objective quantities of
Section~\ref{sec:strain_field}. Writing $\mu = (u_x + v_y)/2$ (zero for
incompressible flow), $a = (u_x - v_y)/2$, and $b = (u_y + v_x)/2$ for
the mean, deviatoric normal, and shear strain, the strain eigenvalues
are $s_{1,2} = \mu \mp r$ with $r = \sqrt{a^2 + b^2}$, and the
eigenframe $(\mathbf{e}_1, \mathbf{e}_2)$ is that of the symmetric
matrix $\hat{\mathbf{S}}$ assembled from the linear coefficients. Under
the quadratic model, $\mu$, $a$, and $b$ are affine functions of
position, so the gradient of $\hat{s}_1$ is exact for the fitted field:
\begin{equation}
    \nabla \hat{s}_1 = \nabla\mu
    - \frac{a\,\nabla a + b\,\nabla b}{r},
    \label{eq:grad_s1}
\end{equation}
with $\nabla\mu$, $\nabla a$, $\nabla b$ read directly off the
second-order coefficients (for example $\nabla a = \tfrac{1}{2}(u_{xx} -
v_{xy},\; u_{xy} - v_{yy})^\top$). Two properties distinguish this
pipeline from the $D$ pipeline, one favorable and one prohibitive.

The favorable one: the tangent the $s_1$ tracker rides,
$\mathbf{e}_2$, comes from the \emph{first-order} fitted coefficients,
whose noise gain is $(2/\sqrt{10})\,\sigma_{uv}/\rho$
(Lemma~\ref{lem:gains}, Section~\ref{sec:sensitivity}), one full
derivative order better than the
$\hat{\mathbf{H}}_D$ eigenvectors the $D$ tracker uses for the
same purpose, which inherit the $(8/\sqrt{10})\,\sigma_{uv}/\rho^2$
gain of the second-order coefficients. Section~\ref{sec:results}
quantifies the difference: at $\sigma_{uv} = 0.01$, the median
$\mathbf{e}_2$ direction error is $4.5^\circ$ against $41^\circ$ for
the $\hat{\mathbf{H}}_D$ eigenvector at the same test point.

\begin{remark}[No usable Hessian of $s_1$ exists under a quadratic fit]
\label{rem:s1hess}
The prohibitive one: the natural transplant of the modified Newton
machinery, a Newton step on $(\nabla \hat{s}_1, \hat{\mathbf{H}}_{s_1})$,
is impossible in principle, not merely inaccurate. Under the quadratic
model the map $\mathbf{p} \mapsto (a, b)$ is affine, so $\hat{s}_1 =
\mu - \|(a, b)\|$ is an affine function minus the norm of an affine
map: a concave function of position, whose Hessian is negative
semidefinite \emph{for every field and every formation}. A trench of
$s_1$ has positive transverse curvature, which the model can therefore
never certify; that curvature lives in third derivatives of the
velocity field, outside the span of any quadratic fit. The situation is
strictly worse than the $\hat{\mathbf{H}}_D$ bias of
Section~\ref{sec:sensitivity}, whose retained second-derivative
products at least preserve the eigen-directions: for the double gyre,
where $b \equiv 0$, the fitted $\hat{\mathbf{H}}_{s_1}$ is identically
zero to numerical precision. The gradient (\ref{eq:grad_s1}) is
unaffected: on the double gyre its transverse component recovers the
trench-restoring signal with a slope of the correct sign and order
(fitted $4.0$ per unit offset against the analytic transverse curvature
$6.9$ at the test point of Section~\ref{sec:results}). The $s_1$ tracker
of Section~\ref{sec:oecs_controller} is designed around exactly this
split: eigenframe and gradient in, Hessian out.
\end{remark}

\subsection{Estimator Accuracy: Noise Gains, Truncation, and Formation
Radius}
\label{sec:sensitivity}

Three error sources separate cleanly. Write $\boldsymbol{\Phi} =
\bar{\boldsymbol{\Phi}}\,\mathbf{S}_\rho$ with $\mathbf{S}_\rho =
\mathrm{diag}(1, \rho, \rho, \rho^2, \rho^2, \rho^2)$, where
$\bar{\boldsymbol{\Phi}}$ depends only on the formation shape. All
scale dependence enters through $\mathbf{S}_\rho^{-1}$.

\subsubsection{Measurement noise}
For the ideal formation the noise amplification has closed form. Let each
measurement carry independent zero-mean noise of standard deviation
$\sigma_{uv}$, as in (\ref{eq:meas_noise}). Since
$\hat{\boldsymbol{\theta}}_u - \boldsymbol{\theta}_u =
\boldsymbol{\Phi}^{-1}\boldsymbol{\eta}$, the standard deviation of each
coefficient is $\sigma_{uv}$ times the Euclidean norm of the
corresponding row of $\boldsymbol{\Phi}^{-1}$.

\begin{lemma}[Exact noise gains of the pentagon-plus-center formation]
\label{lem:gains}
For five robots on a ring of radius $\rho$ and one at the centroid, the
coefficient standard deviations are
\begin{equation}
    \sigma_{u_0} = \sigma_{uv}, \quad
    \sigma_{u_x} = \sigma_{u_y} = \frac{2}{\sqrt{10}}\frac{\sigma_{uv}}{\rho},
    \label{eq:gain_grad}
\end{equation}
\begin{equation}
    \sigma_{u_{xx}} = \sigma_{u_{yy}} = \frac{8}{\sqrt{10}}
    \frac{\sigma_{uv}}{\rho^2},
    \qquad
    \sigma_{u_{xy}} = \frac{4}{\sqrt{10}}\frac{\sigma_{uv}}{\rho^2},
    \label{eq:gain_hess}
\end{equation}
independent of the formation heading, and identically for
$\boldsymbol{\theta}_v$.
\end{lemma}

The proof, by direct inversion of $\boldsymbol{\Phi}$ using the fifth
roots of unity, is in Appendix~\ref{app:proofs}. Two design
consequences follow. First, the gains are isotropic, and the error
covariance of the estimated gradient is
$\tfrac{4}{10}\sigma_{uv}^2\,\rho^{-2}\mathbf{I}$ for every heading:
estimator quality does not degrade as the cluster turns. Second, the
$1/\rho$ and $1/\rho^2$ factors exhibit the fundamental scaling: second
derivatives are harder to sense than first derivatives by one power of
the formation radius, and larger formations average noise better. A
numerical search over six-point formations of the same footprint found
at best a few percent improvement in worst-case Hessian gain, so the
symmetric shape is used without further tuning.

\subsubsection{Truncation bias}
If the field is three times differentiable with
$|\partial^3 u|, |\partial^3 v| \leq M_3$ over the footprint, the Taylor
remainder at robot $i$ is bounded by $\tfrac{M_3}{6}\rho^3$, and passing
it through $\mathbf{S}_\rho^{-1}\bar{\boldsymbol{\Phi}}^{-1}$ bounds the
systematic error by
\begin{equation}
    |\hat{u}_x - u_x| \leq c_1 M_3 \rho^2, \qquad
    |\hat{u}_{xx} - u_{xx}| \leq c_2 M_3 \rho,
    \label{eq:truncation}
\end{equation}
and likewise for the remaining coefficients, where $c_1, c_2$ are
absolute constants of the formation shape (row sums of
$|\bar{\boldsymbol{\Phi}}^{-1}|$ divided by six). The Hessian bias decays
only linearly in $\rho$, matching the $O(\rho)$ bound proved for the
symmetric estimator in \cite{15}: smaller formations track curvature better.

\subsubsection{The radius trade-off}
Combining the two effects, the total error of each fitted second-order
coefficient is bounded by
$c_2 M_3 \rho + c_2' \sigma_{uv}/\rho^2$, which is minimized at
\begin{equation}
    \rho^* = \left(\frac{2 c_2' \sigma_{uv}}{c_2 M_3}\right)^{1/3}
    \propto \left(\frac{\sigma_{uv}}{M_3}\right)^{1/3}.
    \label{eq:rho_star}
\end{equation}
The formation radius is thus a sensing parameter, not merely a spacing
choice: it should grow with the noise floor and shrink with the field's
third-derivative scale. This rule is confirmed for $\hat{D}$ and
$\nabla\hat{D}$ by the formation-radius sweeps reported in
Section~\ref{sec:results}; the derived Hessian $\hat{\mathbf{H}}_D$ is
instead limited by the model bias discussed next.

\subsubsection{Propagation to $D$ and its derivatives}
Errors in the fitted coefficients enter the navigation quantities through
products of coefficient pairs. For two-by-two matrices the expansion of
the determinant is exact at second order:
$\det(\mathbf{J} + \mathbf{E}) = \det(\mathbf{J}) +
\operatorname{tr}(\mathrm{adj}(\mathbf{J})\mathbf{E}) + \det(\mathbf{E})$,
and $\|\mathrm{adj}(\mathbf{J})\|_F = \|\mathbf{J}\|_F$, so the estimated
determinant satisfies
\begin{equation}
    |\hat{D} - D| \;\leq\; \|\mathbf{J}\|_F \|\mathbf{E}\|_F
    + \tfrac{1}{2}\|\mathbf{E}\|_F^2,
    \label{eq:det_perturbation}
\end{equation}
where $\mathbf{E} = \hat{\mathbf{J}} - \mathbf{J}$ collects the four
first-order coefficient errors. The gradient formula (\ref{eq:grad_det})
is complete: differentiating $D = u_x v_y - u_y v_x$ once produces only
products of first- and second-order derivatives, so $\hat{\mathbf{g}}$
inherits errors solely from the fitted coefficients, linearly, with
slopes set by the true derivative magnitudes.

The Hessian is different in kind. Differentiating twice produces terms
such as $u_{xxx} v_y$ and $u_y v_{xxx}$, third derivatives of the field
weighted by the field itself, and a quadratic model has no third
derivatives: (\ref{eq:hess_det}) is the Hessian of the fitted quadratic
surface, not of $D$. The omitted terms constitute a model bias that is
independent of the formation radius and cannot be reduced by averaging;
the situation is that of fitting a line to a curved signal and using the
line's slope where the curvature is unmodeled, acceptable exactly where
the model is locally adequate. For the double gyre at the nominal scale
this bias is large in norm, on the order of $\|\mathbf{H}_D\|_F$ itself
(Section~\ref{sec:results}). Two conclusions nevertheless carry into the
stability analysis: the estimate errors entering the controller are
bounded whenever the coefficient errors are bounded, with explicit
formation-dependent constants; and all such constants degrade
continuously, not abruptly, as the formation deforms, with outright
failure only on the conic configurations of
Proposition~\ref{prop:conic}.

\begin{remark}[Hessian bias]
\label{rem:eigvec}
The controller of Section~\ref{sec:controllers} consumes
$\hat{\mathbf{H}}_D$ only through its eigenvectors and the signs and
rough magnitudes of its eigenvalues, and these degrade far less than
the Frobenius error suggests. $D$ is even in the transverse coordinate
across the separatrix, so both the retained and the omitted terms are
diagonal in the trench frame there: the estimated eigen-directions on
the separatrix are exact to numerical precision, and within a few
degrees at generic points (Section~\ref{sec:results}), while the bias
loads almost entirely onto the eigenvalues. An eigenvalue bias rescales
the transverse Newton snap, which the saturation (\ref{eq:sat}) bounds
and the closed loop absorbs; it does not rotate the commanded
direction. This is why a formation too small to sense curvature
accurately can still track the structure that curvature defines. The
exception is the immediate neighborhood of the wells, where the biased
eigenvalues nearly vanish and their signs become unreliable; the
consequences for terminal capture are examined in
Section~\ref{sec:results}.
\end{remark}

\subsubsection{Empirical accuracy}
\label{sec:est_accuracy}

\begin{figure}[!t]
    \centering
    \includegraphics[width=0.45\textwidth]{figures/estimator_accuracy_vs_noise.png}
    % Generated by experiments/estimator_accuracy_sweep.py (CSV data) and
    % experiments/plot_estimator_accuracy.py (figure), under
    % trunk/Python_Simulations/Vector_Fields/VF_Robot/. Review copy in
    % experiments/outputs/estimator_accuracy/.
    \caption{Relative estimation error of $\hat{D}$, $\nabla\hat{D}$, and
    $\hat{\mathbf{H}}_D$ (median over $2 \times 10^5$ noise draws per point;
    shaded interquartile range in (a)) for the steady double gyre of
    Section~\ref{sec:dg_example}, formation held static at the
    strain-region point $\mathbf{p}_c = (-0.35, -0.20)$. (a) Versus measurement noise $\sigma_{uv}$ at the nominal
    formation radius $\rho = 0.075$; the dashed line marks error equal to
    signal magnitude. (b) Versus formation radius $\rho$ at $\sigma_{uv} =
    0.01$; dotted curves are the noise-free truncation floors. $\hat{D}$
    and $\nabla\hat{D}$ exhibit the radius trade-off of (\ref{eq:rho_star});
    $\hat{\mathbf{H}}_D$ is limited by the structural model bias of
    Section~\ref{sec:sensitivity}.}
    \label{fig:est_accuracy}
\end{figure}

Estimator accuracy was measured open loop: the formation was held static
at the strain-region point $\mathbf{p}_c = (-0.35, -0.20)$ of the steady
double gyre, $2 \times 10^5$ independent noise draws were applied per
condition through the same fit used by the controller, and the estimates
were compared with the analytic values of Appendix~\ref{app:analytic}.
The empirical standard deviation of every fitted coefficient matched the
closed forms of Lemma~\ref{lem:gains} within 0.3\%, and the per-robot
reading error induced by position noise matched the first-order
prediction of Section~\ref{sec:noise_model} within 1.1\%, confirming that
the two noise channels combine into a single effective measurement noise.

Fig.~\ref{fig:est_accuracy} reports the accuracy of the three navigation
quantities as relative error, the error magnitude divided by the
magnitude of the true quantity; a value of one (dashed line) means the
estimate is as wrong as it is large and carries no directional
information. In Fig.~\ref{fig:est_accuracy}(a) the three curves rise
linearly with $\sigma_{uv}$ and sit stacked a constant factor apart: each
additional derivative order multiplies the noise gain by roughly
$1/\rho$, about 13 at the nominal radius, the ladder of
Lemma~\ref{lem:gains}. Reading off the crossings of the dashed line,
$\hat{D}$ remains informative up to $\sigma_{uv} \approx 0.06$ m/s and
$\nabla\hat{D}$ up to $\sigma_{uv} \approx 0.01$ m/s, while
$\hat{\mathbf{H}}_D$ sits at the line even for the quietest sensors:
below $\sigma_{uv} \approx 0.005$ its error is set by the structural
model bias discussed above, not by noise.
Fig.~\ref{fig:est_accuracy}(b) fixes $\sigma_{uv} = 0.01$ and varies the
formation radius. Enlarging the ring averages noise over a longer
baseline, so the solid curves fall; the quadratic model simultaneously
fits the curved field more poorly, so the noise-free truncation floors
(dotted) rise. The compromise of (\ref{eq:rho_star}) appears as minima
at $\rho \approx 0.11$ for $\hat{D}$ and $\rho \approx 0.23$ for
$\nabla\hat{D}$. The optima are properties of this field, location,
and noise level, entering through the third-derivative scale $M_3$;
the one-third-power law is general, but any other field or region
shifts them. $\hat{\mathbf{H}}_D$ exhibits no minimum: its floor is radius
independent, as predicted, and no formation size improves it, though its
eigen-directions remain accurate to $0.35^\circ$ on and near the
separatrix and $3.8^\circ$ at the generic test point
(Remark~\ref{rem:eigvec}).

The gradient threshold locates where closed-loop degradation should
begin: at $\sigma_{uv} \approx 0.01$ a single $\nabla\hat{D}$ estimate
is as wrong as it is right, so any one control decision approaches
chance. The closed loop does not extend this limit: the Monte Carlo
sweeps of Section~\ref{sec:results} place the tracker's 50\% success
level at $\sigma_{uv} \approx 0.01$, coinciding with the
single-estimate threshold. Saturation and the per-cycle redraw of the
noise bound the damage of individual bad commands, but they buy no
additional margin for a traverse that requires a long sequence of
correct decisions.


% ---------------------------------------------------------------------------
\section{Adaptive Navigation Controllers}
\label{sec:controllers}

Table~\ref{tab:trackers} previews the two controllers this section
derives. They share the estimation step, the saturation idiom, and the
three-layer architecture below; every row in which they differ traces
to which term of (\ref{eq:splitting}) the tracker trusts, both terms
for the $D$ tracker's network reach, the strain term alone for
the $s_1$ tracker's frame invariance.
Section~\ref{sec:objectivity_tradeoff} weighs the tradeoff.

\begin{table}[!t]
\centering
\caption{The two trackers at a glance}
\label{tab:trackers}
\footnotesize
\setlength{\tabcolsep}{3pt}
\begin{tabular}{lll}
\hline
 & \textbf{$D$ tracker} & \textbf{$s_1$ tracker} \\
\hline
Surrogate & $D = \det(\mathbf{J})$ & $s_1$ of $\mathbf{S} = (\mathbf{J}+\mathbf{J}^\top)/2$ \\
Invariance & Galilean & objective \\
Structure & connected trench network & bounded attracting segments \\
Tangent from & $\hat{\mathbf{H}}_D$ eigenvectors & $\hat{\mathbf{S}}$ eigenvectors \\
Transverse law & Newton snap & saturated gradient descent \\
Along structure & flow-assisted descent & continuity ride or $s_1$ descent \\
Uses flow $\mathbf{f}$ & FLOW mode and selector & tangent-sign seed only \\
Terminal knob & $D_{\text{capture}}$ & $s_{\text{capture}}$, $s_{\text{trim}}$ \\
Guarantee & network traversal & equivariance, local convergence \\
\hline
\end{tabular}
\end{table}

\subsection{Multilayer Control Architecture}
\label{sec:arch}

The trackers run inside the layered control architecture established in
prior cluster space control studies \cite{14,30}, with three layers. At
the highest level, an adaptive navigation layer ingests the six
synchronous positions and velocity measurements, computes the estimates
of Section~\ref{sec:estimation}, and applies the active tracker's
control law to produce a desired centroid velocity
$\dot{\mathbf{p}}_c^{\,\text{des}}$; the field measurements enter the
architecture only at this layer. A cluster space controller then treats
the formation as a single articulating virtual mechanism: using the
twelve-state cluster parameterization and its analytic inverse
kinematic Jacobian (Appendix~\ref{app:kinematics}), it distributes the
centroid command to six robot velocity commands while regulating the
formation shape to its nominal pentagon-plus-center geometry. At the
lowest level, each robot's controller executes its velocity command
subject to the identified actuator dynamics and speed limits. The lower
two layers operate on velocity commands alone, without reference to the
field; the robot model and all gains are given with the simulation
setup in Section~\ref{sec:robot_model}.

Note on notation: $\alpha_{\text{mom}}$, the momentum coefficient of
the robot model in Section~\ref{sec:robot_model}, is unrelated to
$\alpha_{\text{eig}}$, the real part of a complex eigenvalue pair of
$\mathbf{J}$ used in field classification.

\subsection{The Modified Newton Step}
\label{sec:mod_newton}

Under the quadratic model, $\hat{D}(\mathbf{p})$ is an exact quadratic
surface with gradient $\hat{\mathbf{g}}$ and constant Hessian
$\hat{\mathbf{H}}_D$, so the natural navigation move is the Newton step
$-\hat{\mathbf{H}}_D^{-1}\hat{\mathbf{g}}$, which in the eigenbasis
$\hat{\mathbf{H}}_D \mathbf{w}_k = \lambda_k \mathbf{w}_k$ ($\lambda_1
\leq \lambda_2$ throughout) has per-direction components
$-(\mathbf{w}_k^\top \hat{\mathbf{g}})/\lambda_k$. As a velocity
command the raw step fails twice: its magnitude explodes where
$\hat{\mathbf{H}}_D$ is near singular, and it converges to whatever
stationary point the local model exhibits, deliberately climbing along
a trench wherever the along-trench curvature is negative. Componentwise
saturation through
\begin{equation}
    \mathrm{sat}(c) = v_{\max} \tanh\!\left(\frac{c}{v_{\max}}\right)
    \label{eq:sat}
\end{equation}
repairs the first defect, giving the \emph{attract step}
\begin{equation}
    \dot{\mathbf{p}}_c^{\,\text{des}} = \sum_{k=1}^{2}
    \mathrm{sat}\!\left(-\frac{\mathbf{w}_k^\top
    \hat{\mathbf{g}}}{\lambda_k}\right)\mathbf{w}_k,
    \label{eq:attract_step}
\end{equation}
a bounded constant-speed reach far from the structure that recovers the
exact Newton component, with its superlinear local convergence, once
the argument is small; the one parameter $v_{\max}$ sets both the
cruise speed and the terminal boundary layer. Replacing each eigenvalue
by its magnitude, the saddle-free Newton modification of \cite{35},
repairs the second, giving the \emph{slide step}
\begin{equation}
    \dot{\mathbf{p}}_c^{\,\text{des}} = \sum_{k=1}^{2}
    \mathrm{sat}\!\left(-\frac{\mathbf{w}_k^\top
    \hat{\mathbf{g}}}{|\lambda_k|}\right)\mathbf{w}_k,
    \label{eq:slide_step}
\end{equation}
which agrees with (\ref{eq:attract_step}) in directions of positive
curvature, snapping the centroid onto the trench floor, and reverses
the Newton component into a curvature-normalized descent along the
trench, away from crests, in the direction of negative curvature. Every
term of both laws is quadratic in $\mathbf{w}_k$, so the sign ambiguity
of eigenvectors leaves the commands unchanged.

\subsection{The $D$ Tracker}
\label{sec:sep_controller}

The $D$ tracker composes the two step bodies with a third,
flow-assisted mode. The need for the third mode is visible in the trench
profile of Section~\ref{sec:dg_example}: the slide step descends $D$
along the trench, but at the crest the along-trench gradient vanishes,
and on the upstream side of the crest descent points away from the flow
direction, while the separatrix continues through the crest. On the
separatrix itself, however, the ambient flow is tangent to the structure
and nonzero away from the saddles, so the flow can carry the team across
the crest. The cluster is declared on the separatrix
(FLOW band) when either
\begin{equation}
    |\hat{D}| < \varepsilon_{\text{raw}}
    \quad\text{or}\quad
    \frac{|\hat{D}|}{\|\hat{\mathbf{H}}_D\|_F} < \varepsilon_{\text{dim}},
    \label{eq:flow_band}
\end{equation}
where $\varepsilon_{\text{raw}}$ is a raw threshold and
$\varepsilon_{\text{dim}}$ a curvature-normalized threshold (units of
length squared) that makes the band scale with the local landscape. In
the FLOW band the command follows the measured flow along the trench
tangent $\mathbf{w}_1$ and retains the Newton snap across it: the
\emph{flow step} is
\begin{equation}
    \dot{\mathbf{p}}_c^{\,\text{des}} =
    \mathrm{sat}\!\left(\mathbf{f}^\top \mathbf{w}_1\right)\mathbf{w}_1
    + \mathrm{sat}\!\left(-\frac{\hat{\mathbf{g}}^\top
    \mathbf{w}_2}{\lambda_2}\right)\mathbf{w}_2.
    \label{eq:flow_step}
\end{equation}
Off the band, the local flow direction selects between the attract and
slide laws: if the flow agrees with the along-trench descent direction
the team slides (\ref{eq:slide_step}); if it opposes descent the team
attracts toward the local stationary point (\ref{eq:attract_step}),
which lies at the crest, inside the FLOW band, so the attract law hands
the team over to the flow. In all three modes the tangential motion is
never directed against the ambient flow, a fact used in the stability
analysis. The complete selector, executed once per control cycle after
the estimation step, is:

\vspace{4pt}
\hrule
\vspace{4pt}
\noindent\textbf{$D$ tracker: one control cycle}

\vspace{3pt}
\noindent\textbf{Parameters:} $v_{\max}$, $\varepsilon_{\text{raw}}$,
$\varepsilon_{\text{dim}}$\\
\textbf{Given:} $\hat{D}$, $\hat{\mathbf{g}}$, $\hat{\mathbf{H}}_D$,
$\mathbf{f}$ from (\ref{eq:J_hat})--(\ref{eq:hess_det}), and eigenpairs
$(\lambda_1, \mathbf{w}_1)$, $(\lambda_2, \mathbf{w}_2)$ of
$\hat{\mathbf{H}}_D$, $\lambda_1 \leq \lambda_2$

\vspace{3pt}
\noindent If the band test (\ref{eq:flow_band}) passes (on the
separatrix):\\
\hspace*{1.5em} apply the flow step (\ref{eq:flow_step}).\\
Else, if $\lambda_1 \lambda_2 \geq 0$ ($\hat{D}$ has no saddle; near a
well of $D$):\\
\hspace*{1.5em} apply the attract step (\ref{eq:attract_step}).\\
Else, orient $\mathbf{w}_1$ so that $\hat{\mathbf{g}}^\top \mathbf{w}_1
\leq 0$ (descent along the trench):\\
\hspace*{1.5em} if $\mathbf{f}^\top \mathbf{w}_1 > 0$ (flow agrees with
descent), apply the slide step (\ref{eq:slide_step});\\
\hspace*{1.5em} otherwise (flow opposes descent), apply the attract
step (\ref{eq:attract_step}).
\vspace{4pt}
\hrule
\vspace{4pt}

The implementation is \textit{separatrix\_logic\_c\_step} in
\textit{pentagon\_primitives.py}. One implementation guard is not shown
in the selector: inside the band, if $\hat{\mathbf{H}}_D$ has no saddle
structure, the command falls back to the saturated flow alone. On the
trench with exact estimates the in-band Hessian is a saddle away from
the band boundary, so the guard fires only under noise or off
structure; Section~\ref{sec:results} counts it as the flow-drift
refinement. Note that the $\mathbf{w}_2$
component is one and the same in all three laws
(\ref{eq:attract_step}), (\ref{eq:slide_step}), and
(\ref{eq:flow_step}): on a trench, $\lambda_2 > 0$, so the signed and
magnitude denominators agree, and every mode applies the identical
transverse Newton snap. The modes differ only in their motion along the
structure. This shared component is what makes the switched system
tractable in Section~\ref{sec:stability}.

\subsection{The $s_1$ Tracker}
\label{sec:oecs_controller}

The second controller tracks the attracting OECSs of
Section~\ref{sec:strain_field} using only the objective quantities the
fit delivers: $\hat{s}_1$, $\nabla\hat{s}_1$, and the strain eigenframe
$(\mathbf{e}_1, \mathbf{e}_2)$. Remark~\ref{rem:s1hess} rules out a
Newton snap, so the two channels are built differently from the $D$ tracker's.
The transverse channel descends $\hat{s}_1$ along the compression
eigenvector,
\begin{equation}
    \dot{\mathbf{p}}_c^{\,\perp} =
    \mathrm{sat}\!\bigl(-g_\perp\,\nabla\hat{s}_1^\top
    \mathbf{e}_1\bigr)\,\mathbf{e}_1,
    \label{eq:oecs_perp}
\end{equation}
a saturated gradient step with fixed gain $g_\perp$ in place of the
unavailable curvature normalization; on a trench of $s_1$ the argument
is a restoring signal whose slope is the true transverse curvature, so
(\ref{eq:oecs_perp}) contracts at gradient rate rather than Newton
rate, the first cost of objectivity. The tangential channel drives
along the stretching eigenvector $\mathbf{e}_2$, and a single threshold
$s_{\text{capture}}$ selects between the two behaviors of the
park-versus-traverse pattern (Section~\ref{sec:park_vs_traverse}):
core-seek descends $\hat{s}_1$ along the curve and parks at the TRAP
core, while ride cruises along $\mathbf{e}_2$ at constant speed,
oriented by continuity with the previous step's tangent.
Both channels are quadratic in the eigenvector they use, except the
ride direction, which is exactly where the sign ambiguity of a
direction field cannot be canceled algebraically: the tangent is
seeded from the measured flow once, at the first tracking step, and
thereafter oriented by $\mathrm{sign}(\mathbf{e}_2^\top
\mathbf{t}_{\text{prev}})$ alone, never consulting the flow again. This
is the orientation-by-continuity rule that direction fields with
half-integer singularities require
(Section~\ref{sec:strain_field}); a flow-based rule would surrender
the objectivity the controller exists to provide, since the measured
flow is frame-dependent. The complete selector, executed once per
control cycle after the estimation step, is:

\vspace{4pt}
\hrule
\vspace{4pt}
\noindent\textbf{$s_1$ tracker: one control cycle}

\vspace{3pt}
\noindent\textbf{Parameters:} $v_{\max}$, $g_\perp$,
$s_{\text{trim}}$, $s_{\text{capture}}$ (finite for core-seek, unset
for ride), $\varepsilon_{\text{degen}}$\\
\textbf{Given:} $\hat{s}_1$, $\nabla\hat{s}_1$, $r$, eigenframe
$(\mathbf{e}_1, \mathbf{e}_2)$ from Section~\ref{sec:s1_estimation},
flow $\mathbf{f}$; persistent flags: banded (has $\hat{s}_1 <
-s_{\text{trim}}$ ever held), previous tangent
$\mathbf{t}_{\text{prev}}$

\vspace{3pt}
\noindent If $r < \varepsilon_{\text{degen}}$ (isotropic point, frame
unreliable) or not banded:\\
\hspace*{1.5em} ACQUIRE: apply
$-\mathrm{sat}(g_\perp \partial\hat{s}_1/\partial x),
 -\mathrm{sat}(g_\perp \partial\hat{s}_1/\partial y)$
componentwise.\\
Else, if $\hat{s}_1 > -s_{\text{trim}}$ (attraction no longer
dominates):\\
\hspace*{1.5em} TRIM: hold position (the structure has ended).\\
Else, if $s_{\text{capture}}$ is set and $\hat{s}_1 <
s_{\text{capture}}$ (inside the core well):\\
\hspace*{1.5em} PARK: apply the ACQUIRE step, which settles at the
$s_1$ minimum.\\
Else (TRACK): apply the transverse step (\ref{eq:oecs_perp}), plus\\
\hspace*{1.5em} core-seek: $\mathrm{sat}(-g_\perp\,
\nabla\hat{s}_1^\top \mathbf{e}_2)\,\mathbf{e}_2$ (descend $s_1$ toward
the core);\\
\hspace*{1.5em} ride: $v_{\max}\tanh(1)\,\mathbf{t}$ with $\mathbf{t} =
\pm\mathbf{e}_2$ oriented by $\mathbf{t}_{\text{prev}}$.
\vspace{4pt}
\hrule
\vspace{4pt}

The implementation is \textit{oecs\_trap\_step} in
\textit{pentagon\_primitives.py}. The TRIM test is the Serra-Haller
termination rule playing the role that the eigenvalue capture test
plays in the $D$ tracker, and $s_{\text{capture}}$ is the exact analogue of
$D_{\text{capture}}$: one threshold, two behaviors. Every quantity the
selector consumes transforms equivariantly under a time-dependent
rotation of the observer ($s_1$ and $r$ invariant, $\nabla s_1$ and the
eigenframe co-rotating), so in the tracking branches the commanded
velocity co-rotates as well and the closed-loop trajectory is the same
material path in any frame; the componentwise acquisition steps are
covariant to within their saturation (Section~\ref{sec:stab_oecs}). What the controller gives up relative to
the $D$ tracker is equally explicit: it has no flow-assisted mode, so it does
not traverse crests where $s_1 \to 0$; it tracks attracting structures
only, which on the double gyre are segments, not the full separatrix;
and its transverse convergence is gradient-rate. Section~\ref{sec:results}
shows both sides.

\subsection{Stability Properties}
\label{sec:stability}

This subsection states the closed-loop guarantees for both trackers;
all proofs are collected in Appendix~\ref{app:proofs}.

\subsubsection{Closed-Loop Model and Standing Assumptions}

The analysis is carried out at the navigation level: the cluster
centroid is modeled as a kinematic integrator
\begin{equation}
    \dot{\mathbf{p}}_c = k\,\boldsymbol{\delta}(\mathbf{p}_c),
    \label{eq:kinematic_model}
\end{equation}
where $\boldsymbol{\delta}$ is the output of the $D$ tracker,
and $k > 0$ is the navigation gain. This abstracts the cluster and
robot layers, whose
tracking dynamics (\ref{eq:momentum}) are faster than the navigation
timescale ($\tau \approx 0.3$ s against transit times of tens of seconds);
the same separation was validated in hardware in \cite{14}. Time is rescaled
so that $k = 1$. Unless stated otherwise the estimates are taken exact
($\hat{D} = D$, $\hat{\mathbf{g}} = \mathbf{g}$, $\hat{\mathbf{H}}_D =
\mathbf{H}_D$, $\mathbf{f} = \mathbf{v}(\mathbf{p}_c)$);
Section~\ref{sec:stab_iss} restores bounded estimation error.
Throughout, $\mathbf{w}_1, \mathbf{w}_2$ denote unit eigenvectors of
$\mathbf{H}_D$ with $\lambda_1 \leq \lambda_2$, and $z_k =
(\mathbf{w}_k^\top \mathbf{g})/\lambda_k$ the Newton components, so
that the attract step (\ref{eq:attract_step}) reads
$\boldsymbol{\delta} = -\sum_k \mathrm{sat}(z_k)\,\mathbf{w}_k$.

\subsubsection{The Attract Mode Is a Gradient-Norm Descent}

The first result needs no trench structure at all: it certifies the attract step
globally, wherever $\mathbf{H}_D$ is nonsingular.

\begin{theorem}[Attract-mode certificate]
\label{thm:attract}
Let $V_A = \tfrac{1}{2}\|\mathbf{g}(\mathbf{p}_c)\|^2$. Along
(\ref{eq:kinematic_model}) with the attract step
(\ref{eq:attract_step}),
\begin{equation}
    \dot{V}_A = -v_{\max} \sum_{k=1}^{2} \lambda_k^2\,
    z_k \tanh\!\left(\frac{z_k}{v_{\max}}\right) \;\leq\; 0,
    \label{eq:VA_dot}
\end{equation}
with equality exactly where $\mathbf{g} = \mathbf{0}$. Consequently
every bounded trajectory on which $\mathbf{H}_D$ remains nonsingular
converges to the critical set $\{\mathbf{g} = \mathbf{0}\}$ of $D$, and
isolated critical points with $\mathbf{H}_D$ nonsingular are locally
exponentially stable.
\end{theorem}

Theorem~\ref{thm:attract} explains the two roles the attract step plays
in the selector. Far from the trench it provides the approach: it descends
$\|\nabla D\|^2$ toward the stationary set of $D$, which for the double
gyre consists of the trench line's crest and wells and the gyre
centers. Near a well of $D$, where $\mathbf{H}_D \succ 0$, it is the
terminal capture law: the well is the saddle point of the flow, and
convergence is exponential.

\subsubsection{The $D$ Tracker: Network Traversal}
\label{sec:stab_sep}

The trench is formalized as a valley line. The double gyre satisfies all
of the following exactly, with the constants given below.

\begin{assumption}[Nondegenerate straight trench]
\label{ass:trench}
There is a segment $\Gamma$ (take it along the unit vector
$\hat{\mathbf{e}}_\parallel$, with normal $\hat{\mathbf{e}}_\perp$, and
coordinates $(s, n)$) and a tube $T_w = \{|n| \leq w\}$ on which
\begin{equation}
    D(s, n) = \varphi(s) + \tfrac{1}{2}\lambda_\perp(s)\,n^2,
    \qquad 0 < \underline{\lambda} \leq \lambda_\perp(s),
    \label{eq:normal_form}
\end{equation}
the cross terms vanish
($\hat{\mathbf{e}}_\perp^\top \mathbf{H}_D \hat{\mathbf{e}}_\parallel
= 0$ on $T_w$), and on $T_w$ the eigenvalue gap satisfies
$\lambda_2 - \lambda_1 \geq \gamma > 0$ wherever the selector is
outside its attract fallback. For the double gyre separatrix, take
$n = x$, $s = y$: the field (\ref{eq:det_analytic_main}) is additively
separable, so the cross terms vanish identically and $\varphi(s) =
-\tfrac{\pi^4 A^2}{2}(1 + \cos 2\pi s_f)$ exactly, while the transverse
profile $\pi^4 A^2 \sin^2(\pi n)$ matches the quadratic normal form to
second order in $n$, with $\lambda_\perp = 2\pi^6 A^2$ (the exact
transverse curvature on the trench) and transverse curvature
$2\pi^6 A^2 \cos 2\pi x_f > 0$ throughout any tube with $w < 1/4$; the
proof of Lemma~\ref{lem:transverse} shows that the higher-order terms
only strengthen the contraction.
\end{assumption}

\begin{assumption}[Flow structure on the trench]
\label{ass:flow}
$\Gamma$ is an invariant curve of the flow: $\mathbf{v}$ restricted to
$\Gamma$ is tangent to it, with $\|\mathbf{v}\| \geq f_{\min} > 0$
outside balls of radius $r_0$ around the flow saddles, where the flow
vanishes. Along $\Gamma$, the stationary points of $\varphi$ are of two
kinds only: wells located at the flow saddles, and crests located inside
the FLOW band (\ref{eq:flow_band}); quantitatively, $|\varphi'(s)| \geq
m_\varphi > 0$ wherever the band condition fails on $\Gamma$.
\end{assumption}

\begin{lemma}[Common transverse contraction]
\label{lem:transverse}
Under Assumption~\ref{ass:trench}, in every mode of the $D$ tracker
the transverse coordinate obeys
\begin{equation}
    \dot{n} = -v_{\max}\tanh\!\left(\frac{n}{v_{\max}}\right),
    \label{eq:n_dynamics}
\end{equation}
so $|n|$ decreases strictly for $n \neq 0$, the tube $T_w$ is forward
invariant, and $n \to 0$ exponentially once $|n| \leq v_{\max}$. In
particular (\ref{eq:n_dynamics}) is unaffected by arbitrary switching
among the three modes.
\end{lemma}

\begin{lemma}[Flow-consistent progress]
\label{lem:progress}
Under Assumptions~\ref{ass:trench} and \ref{ass:flow}, on $\Gamma$
(after the transient of Lemma~\ref{lem:transverse}) the tangential
velocity component $\dot{s}$ satisfies, in every mode,
$\dot{s}\,(\mathbf{v}^\top \hat{\mathbf{e}}_\parallel) \geq 0$: the team
never moves against the ambient flow along the trench. Moreover
$|\dot{s}| \geq m > 0$ outside the $r_0$ balls, with
\begin{equation}
    m = \min\Bigl\{
    v_{\max}\tanh\!\bigl(\tfrac{f_{\min}}{v_{\max}}\bigr),\;
    v_{\max}\tanh\!\bigl(\tfrac{m_\varphi}
    {\bar{\lambda}_\parallel v_{\max}}\bigr)
    \Bigr\},
    \label{eq:progress_rate}
\end{equation}
where $\bar{\lambda}_\parallel$ bounds $|\lambda_1|$ on the tube.
\end{lemma}

\begin{theorem}[Separatrix network traversal]
\label{thm:separatrix}
Under Assumptions~\ref{ass:trench} and \ref{ass:flow}, with exact
estimates, every trajectory of the $D$ tracker starting in the
tube $T_w$
(i) remains in $T_w$ and converges exponentially to the trench, and (ii)
traverses the trench in the direction of the ambient flow, reaching the
$r_0$ neighborhood of a flow saddle in finite time bounded by
$L_\Gamma/m$ with $L_\Gamma$ the trench length. At the saddle, $D$
attains a local minimum with $\mathbf{H}_D \succ 0$
($\mathbf{H}_D = 2\pi^6 A^2\,\mathbf{I}$ at the double-gyre saddles), so
the selector's stationary-point test $\lambda_1\lambda_2 \geq 0$
(Section~\ref{sec:sep_controller}) is satisfied there and (iii)
Theorem~\ref{thm:attract} gives local exponential convergence to the
saddle, terminal capture, exactly as stated.
\end{theorem}

\begin{remark}[Role of each mode]
The proof assigns each mode a precise job: the transverse Newton snap
(common to all modes) makes the trench attractive; the slide step produces the
descent along the trench; the flow step carries the team across the crest, where
descent is undefined; the attract step turns upstream starts around by climbing to
the crest, and, when its stationary-point test fires at a well, doubles as the
terminal capture law. The
hybrid has no chattering issue in the transverse channel because all
modes share the dynamics (\ref{eq:n_dynamics}).
\end{remark}

\begin{remark}[Capture depends on an estimated, not exact, Hessian]
\label{rem:capture_gap}
The trench is not a single arc: the saddle is a crossing of two
perpendicular trench segments (Section~\ref{sec:dg_example}), so a
selector that reaches the ball without certifying $\mathbf{H}_D \succ
0$ remains in its flow or slide modes and continues onto the crossing
segment instead, repeating (i)-(ii) with $\Gamma$ relabeled.
Remark~\ref{rem:eigvec} identifies exactly when this happens for the
closed-form estimator: near the wells the biased eigenvalues nearly
vanish and their signs become unreliable, so the stationary-point test
does not fire on $\hat{\mathbf{H}}_D$ even without noise, and the
selector takes continuation rather than capture, which
Section~\ref{sec:results} confirms at every saddle the simulated
controller reaches. Terminal capture is recovered by tightening the
test, demonstrated as a one-parameter alternative in
Section~\ref{sec:discussion}; both behaviors are the same theorem
applied to different premises at the crossing.
\end{remark}

\subsubsection{Closed-Loop Robustness}
\label{sec:stab_iss}

Finally the exact-estimate assumption is dropped. By
Section~\ref{sec:sensitivity}, noise of level $\sigma_{uv}$ and
truncation over a footprint of radius $\rho$ perturb every quantity the
controller uses; the perturbations are bounded, with constants
determined by the formation. Their effect enters each control channel
inside a $\tanh$, which is 1-Lipschitz, so a perturbation of a
commanded component is no larger than the perturbation of its argument.

\begin{proposition}[Ultimate bound under bounded estimate error]
\label{prop:iss}
Suppose the estimate errors induce a bounded disturbance $d(t)$, $|d|
\leq \bar{d} < v_{\max}$, in the transverse channel, so that
(\ref{eq:n_dynamics}) becomes $\dot{n} = -v_{\max}\tanh(n/v_{\max}) +
d$. Then the trench remains practically stable with ultimate bound
\begin{equation}
    \limsup_{t\to\infty} |n(t)| \;\leq\;
    v_{\max}\operatorname{artanh}\!\left(\frac{\bar{d}}{v_{\max}}\right),
    \label{eq:iss_bound}
\end{equation}
which is $\approx \bar{d}$ for $\bar{d} \ll v_{\max}$.
\end{proposition}

The disturbance level $\bar{d}$ inherits the noise and truncation
scalings of Section~\ref{sec:sensitivity}, so the ultimate bound
(\ref{eq:iss_bound}) degrades linearly in the measurement noise level,
with the formation radius trading the two contributions as in
(\ref{eq:rho_star}). The Monte Carlo sweeps of
Section~\ref{sec:test_plan} quantify these predictions empirically.

\subsubsection{The $s_1$ Tracker: Objectivity and Local Convergence}
\label{sec:stab_oecs}

Two properties of the $s_1$ tracker are stated at the same navigation
level as the results above.

\begin{proposition}[Frame equivariance of the closed loop]
\label{prop:equivariance}
Let $\mathbf{p}'= \mathbf{Q}(t)\mathbf{p}$ be a time-dependent rotation
of the observer frame, and let the tracker of
Section~\ref{sec:oecs_controller} run on the transformed field
$\mathbf{v}'(\mathbf{p}', t) = \mathbf{Q}\mathbf{v}(\mathbf{Q}^\top
\mathbf{p}', t) + \dot{\mathbf{Q}}\mathbf{Q}^\top\mathbf{p}'$ with
exact estimates. Then the selector branch chosen at each state is the
same in both frames, and in the TRACK and TRIM branches the commanded
velocity transforms as
$\boldsymbol{\delta}' = \mathbf{Q}\boldsymbol{\delta}$, so the
closed-loop navigation trajectory (\ref{eq:kinematic_model}) satisfies
$\mathbf{p}'_c(t) = \mathbf{Q}(t)\mathbf{p}_c(t)$ up to the
frame-transport term: on the structure, the tracker follows the same
material path in every frame. The ACQUIRE and PARK branches saturate
$\nabla\hat{s}_1$ componentwise in observer axes and are covariant only
where the saturation is inactive; they descend the invariant
$\hat{s}_1$ in every frame, so their contribution is a frame-dependent
transient that the covariant contraction of
(\ref{eq:oecs_n_dynamics}) absorbs.
\end{proposition}

For local convergence, the transverse channel replaces the common
Newton snap of Lemma~\ref{lem:transverse} with the saturated gradient
step (\ref{eq:oecs_perp}). On a trench of $s_1$ satisfying the analogue
of Assumption~\ref{ass:trench}, $s_1(s, n) = \varphi_1(s) +
\tfrac{1}{2}\kappa_\perp(s)\,n^2$ with $0 < \underline{\kappa} \leq
\kappa_\perp$, the transverse coordinate obeys
\begin{equation}
    \dot{n} = -v_{\max}\tanh\!\left(
    \frac{g_\perp \kappa_\perp(s)\,n}{v_{\max}}\right),
    \label{eq:oecs_n_dynamics}
\end{equation}
the same saturated scalar contraction as (\ref{eq:n_dynamics}) with the
unit Newton rate replaced by the field-dependent rate $g_\perp
\kappa_\perp$: exponential once $|n| \leq v_{\max}/(g_\perp
\underline{\kappa})$, with rate proportional to the local curvature
rather than normalized by it. The disturbance analysis of
Proposition~\ref{prop:iss} carries over verbatim with $\bar{d}$ now the
error in $g_\perp \nabla\hat{s}_1^\top\mathbf{e}_1$. Termination is by
the TRIM test rather than by a terminal Lyapunov argument: the tracked
object is a bounded curve segment that ends where $s_1$ loses
negativity, and the selector stops there by construction. In core-seek
mode the PARK step carries its own certificate: it is componentwise
saturated descent of $\hat{s}_1$, so under exact estimates $s_1$ itself
is the Lyapunov function, $\dot{s}_1 = -\sum_k (\partial_k s_1)\,
\mathrm{sat}(g_\perp\,\partial_k s_1) \leq 0$ with equality only at
critical points, and its minimum on the segment is the TRAP core.


% ---------------------------------------------------------------------------
\section{Simulation Studies}
\label{sec:methods}
\label{sec:results}

\subsection{Simulation Setup}

Simulations are run using the Python framework in the
\textit{trunk/Python\_Simulations/Vector\_Fields/VF\_Robot/} directory.
% Entry points: experiments/separatrix_clean_runs.py, ow_clean_runs.py,
% park_vs_traverse_demo.py, mc_sweep_separatrix.py, mc_sweep_ow.py,
% estimator_accuracy_sweep.py, main_ocean_hfr_2km_ftle_overlay.py,
% ocean_hfr_2km_panel_progression.py, ocean_hfr_2km_ftle_snapshots.py,
% ocean_hfr_2km_branch_sensitivity.py (script name comments only, per
% repo style; each script's own header states which paper figure or
% table it produces).
For the double gyre, the field is the static case
(\textit{double\_gyre\_static}, amplitude $A = 0.1$) on the domain
$[-1, 1] \times [-0.5, 0.5]$. The integrator is explicit Euler at
10~Hz ($\Delta t = 0.1$~s), and robot dynamics follow the first-order
model~(\ref{eq:momentum}) with $\alpha_{\text{mom}} = 0.7$. The cluster layer uses
the pentagon-of-pairs formation and its analytic $12 \times 12$ inverse
kinematic Jacobian (Appendix~\ref{app:kinematics}) to distribute the
centroid command to six robot velocity commands, driving each of the
nine formation shape terms to its nominal value proportionally (gain
$1.0$ on lengths, $0.1$ on angles, fixed across all experiments). All
trials use the same nominal formation (\textit{pentagon\_small.yaml})
with ring radius $\rho = 0.075$, the value used throughout
Section~\ref{sec:sensitivity} and Table~\ref{tab:params}; a trial
varies only the formation's initial heading, applied as a rigid
rotation of the base configuration, drawn uniformly from $[0, 2\pi)$ in
the Monte Carlo sweeps and zero in the clean, noise-free demonstration
runs. Clean trials run 400 to 600 steps and stop at domain exit or, for
the $D$ tracker, 150 steps after first saddle contact; noise-sweep
trials run up to 500 steps and stop at task completion or domain exit.
The estimator-accuracy sweep holds the formation static and draws
independent noise per trial rather than stepping the simulator. Every
experiment script under \textit{experiments/} states in a header
comment which paper figure or table it produces and where its CSV or
PNG output is written, so the numbers and figures in this paper trace
back to a single reproducible script each.

\subsection{Robot Model, Gains, and Velocity Limits}
\label{sec:robot_model}

Each robot obeys the first-order momentum model identified on the
Decabot hardware \cite{14},
\begin{equation}
    \mathbf{v}[k+1] = \alpha_{\text{mom}}\,\mathbf{v}[k]
    + (1 - \alpha_{\text{mom}})\,\mathbf{v}^{\text{des}}[k],
    \label{eq:momentum}
\end{equation}
with momentum coefficient $\alpha_{\text{mom}} = 0.7$ at the control
period $\Delta t = 0.1$~s, corresponding to an actuator time constant
$\tau = -\Delta t/\ln\alpha_{\text{mom}} \approx 0.28$~s, consistent
with the $\tau \approx 0.3$~s lag identified on the Decabot hardware;
speed limit $v_{\text{max,robot}} = 0.3$ m/s; and stiction floor
$v_{\text{stiction}} = 0.025$ m/s.
% Code truth: Omnibot default momentum_alpha = 0.7 (omnibot.py). The
% tau = 0.3 s identification would give exp(-1/3) ~= 0.717; the
% simulator uses the fixed 0.7. Do not switch the code value without
% re-running the Monte Carlo baselines (see omnibot.py audit note).
% 0.05 m/s stiction was the three-robot paper's value; the six-robot
% simulation uses the Omnibot default 0.025 (omnibot.py).

The trackers saturate every command component through $v_{\max}\tanh(c /
v_{\max})$, so the commanded centroid speed never exceeds
$\sqrt{2}\,v_{\max}$. The primitive-level saturation speed is $v_{\max}
= 0.04$ m/s, amplified by a fixed navigation gain $k = 3$ before the
cluster layer so that steady commands clear the robots' stiction floor;
the robot-level limits of (\ref{eq:momentum}) are unchanged from \cite{14}.
The band thresholds $\varepsilon_{\text{raw}} = 10^{-3}$ and
$\varepsilon_{\text{dim}} = 0.025$ are fixed in all experiments unless
noted. The $s_1$ tracker shares $v_{\max}$ and $k$ and adds $g_\perp =
1.0$, $s_{\text{trim}} = 0.05$, $\varepsilon_{\text{degen}} = 10^{-3}$,
and, in core-seek mode, $s_{\text{capture}} = -0.9$, all likewise fixed
across experiments ($s_{\text{capture}}$ sits just inside the
double-gyre core wells, whose depth is $s_1 = -\pi^2 A \approx
-0.987$).

Table~\ref{tab:params} collects the values used in all experiments.

\begin{table}[!t]
\centering
\caption{Simulation and controller parameters at the two operating
points: laboratory Decabot (double gyre) and ocean HFR
(Section~\ref{sec:ocean_sim}). Ocean values in world units per
simulated step; one ocean step represents 600~s of field time.}
\label{tab:params}
\footnotesize
\setlength{\tabcolsep}{3pt}
\begin{tabular}{llll}
\hline
\textbf{Symbol} & \textbf{Decabot} & \textbf{Ocean} & \textbf{Description} \\
\hline
$v_{\max}$ & 0.04 & 0.04 & primitive saturation speed (m/s) \\
$k$ & 3.0 & 1.8 & navigation gain \\
$v_{\text{max,robot}}$ & 0.3 & 0.3 & robot speed limit (m/s) \\
$v_{\text{stiction}}$ & 0.025 & 0.002 & stiction floor (m/s) \\
$\tau$ & 0.28 & -- & actuator time constant $-\Delta t/\ln\alpha_{\text{mom}}$ (s) \\
$\Delta t$ & 0.1 & 0.1 (= 600 s field) & control period (s) \\
$\alpha_{\text{mom}}$ & 0.7 & 0 & momentum coefficient \\
$A$ & 0.1 & -- & double-gyre amplitude \\
$\rho$ & 0.075 & 0.105 ($1.4\times$) & formation ring radius (m) \\
$\varepsilon_{\text{raw}}$ & $10^{-3}$ & $10^{-3}$ & FLOW-band raw threshold (s$^{-2}$) \\
$\varepsilon_{\text{dim}}$ & 0.025 & 0.025 & FLOW-band dimensionless thr. (m$^2$) \\
$g_\perp$ & 1.0 & 1.0 & TRAP $s_1$-descent gain (m$^2$) \\
$s_{\text{trim}}$ & 0.05 & 0.3 & TRAP attraction-dominance thr. (s$^{-1}$) \\
$s_{\text{capture}}$ & $-0.9$ & $-3.0$ & TRAP core-park threshold (s$^{-1}$) \\
$\varepsilon_{\text{degen}}$ & $10^{-3}$ & $10^{-3}$ & eigenframe degeneracy thr. (s$^{-1}$) \\
\hline
\end{tabular}
\end{table}

\subsection{Ocean HFR Setup}
\label{sec:ocean_sim}

% --- DATA PROVENANCE AND REPRODUCTION INSTRUCTIONS ---
% The 29 NetCDF files used here are the USWC Radial-to-Total Vector (RTV)
% product at 6 km resolution, produced by CORDC/Scripps and archived at NOAA
% NCEI under accession gov.noaa.nodc:IOOS-HFRadarRTVector.
%
% To re-download the exact 29 hourly files used in this experiment:
%
%   BASE=https://www.ncei.noaa.gov/thredds-ocean/fileServer/ioos/hfradar/rtv/2012/201205/USWC
%
%   for HOUR in 0800 0900 1000 1100 1200 1300 1400 1500 1600 1700 \
%               1800 1900 2000 2100 2200 2300 0000 0100 0200 0300 \
%               0400 0500 0600 0700 0800 0900 1000 1100 1200; do
%     # May 16 files use date prefix 20120516; May 17 files use 20120517.
%     # Hours 0800-2300 on May 16 are prefix 201205160800 through 201205162300.
%     # Hours 0000-1200 on May 17 are prefix 201205170000 through 201205171200.
%     curl -O "$BASE/${TIMESTAMP}_hfr_uswc_6km_rtv_uwls_SIO.nc"
%   done
%
% Construct TIMESTAMP as YYYYMMDDHHOO where YYYYMMDD is the calendar date and
% HHOO is the UTC hour plus two zero digits (e.g., 201205160800 for
% May 16 08:00 UTC). The server returns HTTP 200 directly; no authentication
% or browser session is required.  The files were originally batch-exported
% from the CORDC internal archive on 2020-01-21 by user "motero" using the
% tool rtvUpdateNetcdf (visible in each file's NC history attribute).
% The UCSD/Scripps THREDDS server (hfrnet-tds.ucsd.edu) was decommissioned
% June 30, 2025; do not attempt to use it.
%
% A 2 km resolution version of the same 29 hours is available at the same
% NCEI path with the token "2km" in place of "6km".  See
% trunk/Python_Simulations/Vector_Fields/ocean_data/DATA_PROVENANCE.md
% for a full inventory and notes on spatial coverage differences.
% --- END DATA PROVENANCE ---

The ocean experiment uses hourly surface-current frames from the
Santa Barbara Channel HF radar archive (May 16--17, 2012, 29 hourly
frames covering 28~h of real time). The data are the U.S.~West~Coast
Radial-to-Total Vector (RTV) product produced by the Scripps Coastal
Ocean Observing Laboratory (CORDC) HFRNet network and archived at the
NOAA National Centers for Environmental Information under accession
\textit{gov.noaa.nodc:IOOS-HFRadarRTVector}, retrieved from the NOAA
THREDDS-Ocean server in June 2026. The $2$~km resolution files were
used; each file contains one hour of surface-current velocity
components ($u$, $v$) on a regular geographic grid. Frames are
linearly interpolated in time inside the field evaluator
(\textit{ocean\_hfr\_socal\_timevarying} in
\textit{src/fields/environments/Ocean\_HFR.py}). The cluster runs in
non-dimensional world coordinates on $[-0.65, 0.65]^2$, and an affine
map carries those coordinates to a geographic region of radius
$0.3^{\circ}$ centered on $(34.2^{\circ}\mathrm{N},
-120.4^{\circ}\mathrm{W})$, the same region as the $6$~km product; one
world unit corresponds to approximately $51$~km of latitude at this
center. The pentagon formation is enlarged to $1.4\times$ the nominal
ring radius, $\rho = 0.105$ world units
(\textit{pentagon\_small\_2km.yaml}): the $2$~km product's
lower valid-data coverage ($48$--$58\%$ versus $67\%$ at $6$~km, see
\textit{ocean\_data/DATA\_PROVENANCE.md}) makes it noisier than the
$6$~km field, and at the nominal radius $\rho = 0.075$ the formation
was observed to pick the wrong branch at the bifurcation near
$(34.1^{\circ}\mathrm{N}, -120.4^{\circ}\mathrm{W})$; averaging over a
larger footprint recovers the branch the $6$~km product tracks.

The simulator's discrete-time step is treated as one sample per ten
minutes of real time (\textit{TIME\_WARP} $= 6000$), rather than the
testbed's native $0.1$~s; the Ocean column of Table~\ref{tab:params}
collects the resulting operating point alongside the laboratory
Decabot values. Two changes follow directly from the time dilation:
the momentum coefficient is set to $\alpha_{\text{mom}} = 0$, since at
one control cycle per ten real-world minutes any physically reasonable
vessel time constant is negligible, and the stiction floor is lowered
to $0.002$, since the Decabot floor models motor static friction,
which has no vessel-scale analogue and would incorrectly clip small
commanded corrections. The navigation gain is retuned from $k = 3.0$
to $k = 1.8$, with $v_{\max} = 0.04$; empirically this lets the
tracker follow the ridge branch that descends along the middle
island's north coast and makes landfall there, rather than continuing
east past it. With $168$ sample steps the trial spans $168 \cdot
10$~min $= 28$~h of real field time, the full duration of the archived
record.

\subsection{Test Plan}
\label{sec:test_plan}

Four experiment families validate the trackers.

\subsubsection{Double-gyre clean runs}
Each tracker runs noise-free from six initial conditions spanning both
gyre halves, the crest, and a near-center weak-strain point; a trial
stops at domain exit or a fixed number of steps after task completion.
These runs verify the designed mode sequences and terminal behaviors
against the analytic geometry of Section~\ref{sec:dg_example}. A
park-versus-traverse pair (Section~\ref{sec:park_vs_traverse}) runs
both settings of the terminal threshold from the same start,
noise-free.

\subsubsection{Monte Carlo noise sweeps}
Noise robustness is evaluated over a two-dimensional noise grid, from
one consistent initial condition, a formation straddling the
separatrix at $(0, 0.35)$, so that the success rate measures noise
robustness alone rather than basin geometry
(Section~\ref{sec:limitations}). The grid spans $\sigma_{uv} \in \{0,\,
0.001,\, 0.005,\, 0.01,\, 0.02,\, 0.05,\, 0.1\}$ m/s, with the top
level doubled and re-swept, up to three times ($\sigma_{uv} = 0.8$),
targeting a worst-cell success rate below 10\% so the failure cliff is
bracketed (the $D$ tracker reaches that level; the $s_1$ tracker's
random-walk floor does not), and $\sigma_p \in
\{0,\, 0.005,\, 0.01,\, 0.02,\, 0.05\}$ m, at 1000 trials per cell of
at most 500 steps each. Every trial draws an independent initial
heading, testing the isotropy prediction of Lemma~\ref{lem:gains}.
Task success is reaching the tracker's terminal target (centroid
within 0.06) without formation collapse (root-mean-square error of the
fifteen inter-robot distances from their nominal values exceeding
0.30, four times the formation radius $\rho$); each trial also
records time-to-band, straddle retention (whether robots remain on
both sides of the separatrix from first straddle to stop, the failure
count of \cite{9}), tracking error (mean and 95th percentile of
$|x_c|$ over the on-structure phase), and control effort.

The terminal targets necessarily differ. The $D$ tracker traverses to
the terminal saddle $\mathbf{p}^*_1$, while the attracting segment
through the same start ends at its TRAP core, the top saddle
$\mathbf{p}^*_2 = (0, 0.5)$, so the $s_1$ tracker, swept in core-seek
mode with the same random-heading draw, is scored on core arrival. The
on-segment path is shorter than the separatrix traverse (0.15 against
0.85), so success rates compare cell by cell only qualitatively; the
prediction under test is the location of the failure cliff, which the
estimator analysis places at the same $\sigma_{uv} \approx 0.01$ for
both controllers, since $\nabla\hat{s}_1$ and $\nabla\hat{D}$ share
the second-order coefficient noise gain, while straddle retention and
tracking error test whether the first-order eigenframe
(Section~\ref{sec:s1_estimation}) buys measurable robustness on the
structure.

\subsubsection{Objectivity demonstration}
Both controllers run on the same physical double gyre expressed in two
observer frames: the inertial frame, and a frame rotating at
$\Omega = 0.2$ rad/s, in which the field is
$\mathbf{v}'(\mathbf{p}', t) = \mathbf{Q}\mathbf{v}(\mathbf{Q}^\top
\mathbf{p}') + \Omega\,(-y', x')^\top$. The added solid-body vorticity
$2\Omega = 0.4$ is 20\% of the gyre-core vorticity and dominates the
near-zero vorticity at the separatrix; the rate was chosen so the
structure's apparent speed at the tracked radius, $\Omega r \leq
0.10$, stays inside the command authority $k\,v_{\max} = 0.12$.
Rotating-frame trajectories are pulled back to inertial coordinates
through $\mathbf{Q}^\top(t)$ and compared with the inertial run of the
same controller; Proposition~\ref{prop:equivariance} predicts
near-zero gaps for the $s_1$ tracker and offers no such protection to
the $D$ tracker, whose surrogate obeys (\ref{eq:D_not_objective}).

\subsubsection{Ocean runs}
Both trackers run unmodified on the Santa Barbara Channel field of
Section~\ref{sec:ocean_sim}, at the Ocean operating point of
Table~\ref{tab:params}, and are evaluated against ground truth
computed offline from the full radar grid: 24-h forward-time FTLE
fields for the $D$ tracker and TRAP cores (local minima of $s_1$ per
hourly frame) for the $s_1$ tracker. A $10 \times 10$ grid of jittered
starts tests the $D$ tracker's sensitivity to its initial condition.

\subsection{Double Gyre: The $D$ Tracker}
\label{sec:results_sep}

\begin{figure}[!t]
    \centering
    \includegraphics[width=0.45\textwidth]{figures/separatrix_trajectories.png}
    % Generated by experiments/separatrix_clean_runs.py, under
    % trunk/Python_Simulations/Vector_Fields/VF_Robot/. Review copy in
    % experiments/outputs/separatrix_clean/.
    \caption{$D$ tracker centroid trajectories from six
    initial conditions, noise-free case, overlaid on the double-gyre
    streamlines and the two Okubo-Weiss diamonds. Start points are
    $\bullet$, saddles are $\times$. Every trajectory reaches $x = 0$ and
    continues along the wall trench through the terminal (bottom) saddle
    in both directions; a trial stops at domain exit or 150 steps after
    first saddle contact.}
    \label{fig:sep_trajectories}
\end{figure}

From the six starts of Fig.~\ref{fig:sep_trajectories}, all six reach the
separatrix and pass within 0.02 of the bottom saddle
$\mathbf{p}^*_1$ (closest approach 0.0003--0.0194 across runs). Two of
the six begin already inside the FLOW band and band immediately;
the other four take 6, 13, 18, and 35 steps. Mode occupancy over each run
is FLOW 6--37\%, SLIDE 34--48\%, and ATTRACT 10--24\%, with the
remaining time split between the flow-drift and attract-fallback
refinements of the same three modes; no run spends more than a few
percent of its steps outside these categories. Every run continues past
the terminal saddle rather than stopping there. This traces to the
structural bias of $\hat{\mathbf{H}}_D$ (Section~\ref{sec:sensitivity}):
near the well the estimated eigenvalues are indefinite and more than an
order of magnitude too small, so the capture test
$\lambda_1\lambda_2 \geq 0$ of Section~\ref{sec:sep_controller} does not
fire even without noise. Section~\ref{sec:discussion} returns to this
behavior, continuous traversal of the separatrix network, with terminal
capture demonstrated as a tightened-threshold special case.

\begin{table*}[!t]
\centering
\caption{Monte Carlo noise robustness of both trackers: 1000
trials/cell, fixed straddling start $(0, 0.35)$, random heading per
trial. (a), (b): task success (\%), traverse to the terminal saddle
$\mathbf{p}^*_1$ for the $D$ tracker, core arrival at
$\mathbf{p}^*_2$ for the $s_1$ tracker. (c), (d): straddle retention
(fraction of trials, the Michini-comparable metric \cite{9}). Both
sweeps auto-extended to $\sigma_{uv} = 0.8$; rows beyond 0.2 are
omitted: there the $s_1$ tracker sits on its 34--38\% random floor,
while the $D$ tracker decays from 11--13\% at 0.4 to 5--6\% at 0.8.}
\label{tab:mc_success}
\footnotesize
\setlength{\tabcolsep}{4pt}
\begin{tabular}{r|rrrrr}
\multicolumn{6}{l}{(a) $D$ tracker: traverse success (\%)} \\
\hline
$\sigma_{uv}$ \textbackslash\ $\sigma_p$ & 0.000 & 0.005 & 0.010 & 0.020 & 0.050 \\
\hline
0.000 & 100.0 & 64.7 & 45.4 & 36.0 & 18.8 \\
0.001 &  99.3 & 61.9 & 43.7 & 33.8 & 19.9 \\
0.005 &  60.2 & 53.6 & 40.7 & 29.7 & 17.4 \\
0.010 &  43.3 & 40.6 & 39.6 & 28.1 & 18.4 \\
0.020 &  27.6 & 27.5 & 22.9 & 20.4 & 18.4 \\
0.050 &  15.1 & 16.4 & 14.7 & 15.2 & 15.0 \\
0.100 &  13.0 & 15.0 & 13.3 & 14.3 & 13.5 \\
0.200 &  13.9 & 11.8 & 15.4 & 13.6 & 11.0 \\
\hline
\end{tabular}
\hspace{1.5em}
\begin{tabular}{r|rrrrr}
\multicolumn{6}{l}{(b) $s_1$ tracker: core-arrival success (\%)} \\
\hline
$\sigma_{uv}$ \textbackslash\ $\sigma_p$ & 0.000 & 0.005 & 0.010 & 0.020 & 0.050 \\
\hline
0.000 & 100.0 & 98.6 & 94.3 & 87.6 & 79.6 \\
0.001 & 100.0 & 98.4 & 94.6 & 88.8 & 81.2 \\
0.005 &  97.7 & 95.8 & 93.5 & 87.1 & 80.0 \\
0.010 &  94.3 & 92.7 & 89.8 & 85.4 & 80.4 \\
0.020 &  84.6 & 86.0 & 84.9 & 80.5 & 75.0 \\
0.050 &  70.6 & 70.2 & 67.8 & 68.5 & 65.3 \\
0.100 &  51.9 & 48.1 & 53.9 & 51.3 & 50.4 \\
0.200 &  40.1 & 40.9 & 37.7 & 42.6 & 42.6 \\
\hline
\end{tabular}
\\[8pt]
\begin{tabular}{r|rrrrr}
\multicolumn{6}{l}{(c) $D$ tracker: straddle retention} \\
\hline
$\sigma_{uv}$ \textbackslash\ $\sigma_p$ & 0.000 & 0.005 & 0.010 & 0.020 & 0.050 \\
\hline
0.000 & 1.000 & 0.625 & 0.417 & 0.338 & 0.164 \\
0.005 & 0.560 & 0.493 & 0.365 & 0.272 & 0.157 \\
0.010 & 0.379 & 0.366 & 0.361 & 0.258 & 0.165 \\
0.020 & 0.237 & 0.232 & 0.187 & 0.177 & 0.159 \\
0.050 & 0.126 & 0.137 & 0.112 & 0.128 & 0.122 \\
\hline
\end{tabular}
\hspace{1.5em}
\begin{tabular}{r|rrrrr}
\multicolumn{6}{l}{(d) $s_1$ tracker: straddle retention} \\
\hline
$\sigma_{uv}$ \textbackslash\ $\sigma_p$ & 0.000 & 0.005 & 0.010 & 0.020 & 0.050 \\
\hline
0.000 & 1.000 & 0.922 & 0.749 & 0.535 & 0.334 \\
0.005 & 0.894 & 0.829 & 0.722 & 0.511 & 0.356 \\
0.010 & 0.719 & 0.673 & 0.625 & 0.481 & 0.345 \\
0.020 & 0.483 & 0.512 & 0.441 & 0.402 & 0.319 \\
0.050 & 0.271 & 0.281 & 0.284 & 0.268 & 0.225 \\
0.100 & 0.189 & 0.186 & 0.199 & 0.174 & 0.167 \\
0.200 & 0.137 & 0.163 & 0.138 & 0.179 & 0.154 \\
\hline
\end{tabular}
\end{table*}

% Sweep implementation: experiments/mc_sweep_separatrix.py (VF_Robot);
% analysis by experiments/analyze_mc_sweeps.py; per-trial rows in
% experiments/outputs/mc_*/trials_fixed.csv.
Table~\ref{tab:mc_success}(a) reports traverse success (reaching the
terminal saddle without formation collapse) over the noise grid. At
$\sigma_p = 0$ the rate is a clean step: 100\% at zero noise and
99.3\% at $\sigma_{uv} = 0.001$, then falling through the 50\% level between
$\sigma_{uv} = 0.005$ and $0.01$, landing at 43.3\% at
$\sigma_{uv} = 0.01$, coinciding with the open-loop $\nabla\hat{D}$
reliability threshold reported above: the estimator, not the control
law, sets the noise ceiling. Position noise degrades the rate along the same curve scaled
by $\|\mathbf{J}\|$, consistent with the effective-noise argument of
Section~\ref{sec:noise_model}: $\sigma_p = 0.01$ alone gives 45.4\%,
close to $\sigma_{uv} = 0.01$ alone at 43.3\%. Beyond the cliff, success
does not fall to zero but levels off near 11--15\% through
$\sigma_{uv} = 0.2$, decaying to 5--6\% only at the extreme extension
level $\sigma_{uv} = 0.8$: at this noise level
the controller behaves as an unbiased random walk that occasionally
lands on the saddle by chance, so the straddle-retention metric of
Table~\ref{tab:mc_success}(c) is the sharper failure indicator. Tracking
error over the on-trench phase (mean $|x_c|$) rises from 0.002 m at zero
noise to 0.086 m at the $\sigma_{uv} = 0.01$ cliff and saturates near
0.27 m at high noise (comparable to the domain half-width), reflecting
trials that never band at all. Success was independent of initial
heading to within the sampling noise of the sweep (quartile success
74.9--76.2\% at low noise, confirming the isotropy of
Lemma~\ref{lem:gains}).

\subsection{Double Gyre: The $s_1$ Tracker}
\label{sec:results_oecs}

% Clean runs: experiments/main_separatrix_oecs.py
% (outputs/oecs/oecs_clean_runs.csv); estimator diagnosis:
% experiments/oecs_estimator_check.py (outputs/oecs/
% oecs_estimator_check.txt, _noise.csv); objectivity:
% experiments/oecs_objectivity_demo.py (figures/oecs_objectivity.png).

\subsubsection{Noise-free runs}
From six noise-free starts spanning both gyre halves, the crest, and a
near-center weak-strain point, the core-seek tracker reached and parked
at the TRAP core predicted by the segment geometry of
Section~\ref{sec:dg_example} in all six cases: the three upper-half
starts and the near-center start at the top saddle, the crest start and
the lower-half start at the bottom saddle, with final distances
0.011--0.014 from the core and zero position drift over the final 50
steps. Time-to-band was 0--3 steps everywhere, including the crest
start, which sits on a degeneracy line of $\mathbf{S}$ and must first
escape it through the ACQUIRE mode. In ride mode the same starts ride
their attracting segment and stop where the theory says the structure
ends: the upper-half starts descend the separatrix and hold at the
crest degeneracy (final $y$ within 0.006 of the isotropic line), the
lower-half starts ride the bottom wall and hold at the wall-corner
degeneracy near $(0.5, -0.5)$. The contrast with
Fig.~\ref{fig:sep_trajectories} is the designed one: the $D$ tracker's flow
mode carries it through the crest and its trench continuation carries
it through the terminal saddle, while the $s_1$ tracker treats both as
what they are objectively, the end of one attracting structure and the
core of another.

\subsubsection{Objectivity}
% Implementation: experiments/oecs_objectivity_demo.py;
% figure figures/oecs_objectivity.png, CSV outputs/oecs/.
The objectivity demonstration follows the rotating-frame protocol of
Section~\ref{sec:test_plan}. Fig.~\ref{fig:oecs_objectivity} shows the
outcome. The
$s_1$ tracker's two runs are the same material path: final gap 0.020,
mean gap 0.017 over the common steps, parking 0.011 (inertial) and
0.009 (rotating) from the top saddle, with mean distance to the true
trench network of 0.001 and 0.005 respectively. The $D$
tracker's two runs are not: the rotating observer's $D'$ landscape (\ref{eq:D_not_objective})
and swirl-contaminated flow measurement carry the team off the
separatrix entirely, to a mean trench-network distance of 0.052 (ten
times its inertial value of 0.005) and a final gap of 1.219 from its
inertial twin, ending in the gyre interior where no structure exists.
The figure demonstrates rather than asserts the claim of
Section~\ref{sec:strain_field}: under a rotating observer the
$D$-based tracker follows a frame artifact, and the $\mathbf{S}$-based
tracker does not.

\begin{figure}[!t]
    \centering
    \includegraphics[width=0.48\textwidth]{figures/oecs_objectivity.png}
    % Generated by experiments/oecs_objectivity_demo.py, under
    % trunk/Python_Simulations/Vector_Fields/VF_Robot/. Review copy in
    % experiments/outputs/oecs/.
    \caption{Objectivity experiment. Each controller runs on the same
    physical double gyre in the inertial frame (solid) and in a frame
    rotating at $\Omega = 0.2$ rad/s (dashed, pulled back to inertial
    coordinates through $\mathbf{Q}^\top(t)$), from the same start
    ($\bullet$). Left: the $s_1$ tracker's two trajectories coincide
    (final gap 0.020) and park at the TRAP core ($\times$). Right:
    the $D$ tracker's rotating-frame run leaves the separatrix and ends off
    structure (final gap 1.219), the frame artifact predicted by
    (\ref{eq:D_not_objective}).}
    \label{fig:oecs_objectivity}
\end{figure}

\subsubsection{Noise robustness}
% Sweep: experiments/mc_sweep_oecs.py, 1000 trials/cell, fixed start
% (0, 0.35), core-seek; table from outputs/mc_oecs/summary_fixed.csv,
% git 17a2c72. Auto-extension ran to sigma_uv = 0.8; the 0.4 and 0.8
% rows sit on the random floor (34-38%) and are summarized in prose.
Table~\ref{tab:mc_success}(b) reports core-arrival success for the
$s_1$ tracker over the same noise grid, with straddle retention in
panel (d). The zero-noise tracking error is 0.008, four times the $D$
tracker's 0.002, the gradient-rate cost predicted by
Remark~\ref{rem:s1hess}. Everything else moves in the $s_1$ tracker's
favor. At $\sigma_{uv} = 0.01$, the level that halves the separatrix
traverse, core arrival still succeeds in 94.3\% of trials; the 50\%
level is not reached until $\sigma_{uv} \approx 0.1$, and beyond it
success settles onto a floor of 34--38\% (measured out to $\sigma_{uv}
= 0.8$) rather than collapsing, since the short on-segment path lets a
noise-driven walk reach the 0.06 contact disk by chance. Straddle
retention is free of that floor and places the sharp degradation
between $\sigma_{uv} = 0.01$ and $0.02$, still two to four times the
$D$ tracker's cliff (panels (d) against (c)). Position noise is
where the first-order eigenframe shows most clearly: $\sigma_p = 0.05$
alone, which leaves the $D$ tracker at 18.8\%, leaves the $s_1$
tracker at 79.6\%. Part of the margin is the shorter task, inherent to
the objective geometry; the rest is visible in the tracking-error
columns, which degrade smoothly for the $s_1$ tracker but jump for the
$D$ tracker. The
mechanism is the one Section~\ref{sec:s1_estimation} predicts: the
tangent direction survives noise an order of magnitude past the point
where every gradient-based quantity approaches chance, so a noisy team
keeps riding the right line even while its transverse regulation
coarsens.



\subsection{Ocean HFR: The $D$ Tracker}
\label{sec:results_ocean}

The $D$ tracker was run on this field from a single start north
of the Santa Barbara Channel entrance, $(34.40^{\circ}\mathrm{N},
-120.39^{\circ}\mathrm{W})$, using the operating point of
Table~\ref{tab:params}. Over the 28~h trial the cluster descends
through the channel entrance, tracks the ridge through the bifurcation
near $(34.1^{\circ}\mathrm{N}, -120.4^{\circ}\mathrm{W})$, and
continues to landfall on the middle island's north coast.
% Progression figure available: figures/ocean_panel_progression_2km.png
% (experiments/ocean_hfr_2km_panel_progression.py); cut for space, the
% full path is visible in the FTLE overlay figure.
Fig.~\ref{fig:ocean_ftle}
places the trial against the 24-h forward-time FTLE field, computed
offline over the same window in the style of \cite{17}; the cluster's path
follows the dominant FTLE ridge running from the channel entrance to the
bifurcation closely, without the ridge being available to the controller
at any point, only the local $\hat{D}$ and $\hat{\mathbf{g}}$ estimates
that Section~\ref{sec:estimation} builds from the current samples at the
formation.

\begin{figure}[!t]
    \centering
    \includegraphics[width=0.48\textwidth]{figures/ocean_ftle_trajectory_overlay_2km.png}
    % Generated by experiments/main_ocean_hfr_2km_ftle_overlay.py, under
    % trunk/Python_Simulations/Vector_Fields/VF_Robot/. Full-resolution
    % copy in ocean_data/det_jacobian_plots/ftle_trajectory_overlay_2km.png.
    \caption{Full 28-h centroid and robot paths overlaid on the 24-h
    forward-time FTLE field computed from the same 2~km data. The path
    tracks the dominant ridge from the channel entrance to the
    bifurcation and continues to landfall near
    $(34.04^{\circ}\mathrm{N}, -120.26^{\circ}\mathrm{W})$.}
    \label{fig:ocean_ftle}
\end{figure}

Forward FTLE recomputed at four anchor times across the record (0, 5,
10, and 16~h, each integrated 12~h forward) confirms that the dominant
ridge through the channel entrance persists while the finer structure
downstream evolves: the experiment is genuinely time-varying, the
tracked structure is not fixed in advance, and the controller sees only
the instantaneous samples at its own formation position.
% Snapshot figure available: figures/ocean_ftle_snapshots_2km.png
% (experiments/ocean_hfr_2km_ftle_snapshots.py); cut from 2A for space.

Fig.~\ref{fig:ocean_branch} tests the trial's sensitivity to the start
position with a $10 \times 10$ grid of starts jittered
$\pm 0.05^{\circ}$ around the validated start, each run through the
full 28-h trial. Of the 100 starts, 84 land on the same
middle-island branch as the validated trial; the remaining 16, all on
or within one grid cell of the jitter-box edges furthest from the
ridge, are carried onto a
western branch that does not make landfall in the domain. The mean
distance from a start to the 95th-percentile FTLE ridge contour was
1.4~km, small relative to the $\sim 0.3^{\circ}$ ($\sim 33$~km) domain
radius, so the branch boundary sits close to the jitter box: most of
the box is on one side of it. This is a basin-of-attraction result in
the sense of Section~\ref{sec:limitations}, not a robustness guarantee;
it shows that the tracked branch is the typical outcome for starts near
the validated one, not that every start converges to it.

\begin{figure}[!t]
    \centering
    \includegraphics[width=0.48\textwidth]{figures/ocean_branch_sensitivity_2km.png}
    % Generated by experiments/ocean_hfr_2km_branch_sensitivity.py, under
    % trunk/Python_Simulations/Vector_Fields/VF_Robot/. Per-trial data in
    % ocean_data/det_jacobian_plots/ocean_branch_sensitivity_2km.csv.
    \caption{Landfall branch for a $10 \times 10$ grid of starts jittered
    around the validated start, over the 24-h forward FTLE background.
    84 of 100 starts (blue) reach the same middle-island branch as
    Fig.~\ref{fig:ocean_ftle}; 16 (red), all near the jitter box edges,
    are carried onto a western branch instead.}
    \label{fig:ocean_branch}
\end{figure}

\subsection{Ocean HFR: The $s_1$ Tracker}
% Run: experiments/main_ocean_hfr_2km_oecs.py; figure
% figures/ocean_oecs_overlay_2km.png; CSV outputs/oecs/
% ocean_oecs_summary.csv. Thresholds from the probed s1 scale of the
% 2km field (median -0.73, p5 -3.0 over the domain at t=0).
The core-seek tracker was run on the real 2~km field at the operating
point of Table~\ref{tab:params}, from the same start as the
separatrix run, with the strain thresholds set from the field's own
$s_1$ scale ($s_{\text{trim}} = 0.3$, $s_{\text{capture}} = -3.0$
against a domain median of $-0.73$). Ground truth is computed offline
from the full radar grid in the style of \cite{17,36}, TRAP cores as
local minima of $s_1$ at each hourly frame; unlike the double gyre,
the ocean strain field is generic, with isolated degeneracies and
cores that move between frames, so this experiment also exercises the
continuity-orientation rule on a live field.
Fig.~\ref{fig:ocean_oecs} shows
the outcome: the team converges to a local minimum of
$s_1$ near the channel entrance within the first hours and
station-keeps on it for the remainder of the 28~h record, ending
0.8~km, less than one grid cell of the 2~km product, from the nearest
TRAP core extracted offline from the full radar grid at the final
frame (local minima of $s_1$ below the 15th percentile, $s_1 <
-1.9$). The contrast with the separatrix run of
Fig.~\ref{fig:ocean_ftle} is the same one the double gyre showed in
miniature: the $D$ tracker rides a ridge sixty kilometers to landfall, while
the objective tracker treats the strongest local attraction point as
its terminal object and holds it as the field evolves. The formation
footprint sets an observable scale limit: the fitted $\hat{s}_1$ at
the parked position averages $-1.4$ over the formation footprint (ring
radius 0.105 world units, $\sim$5.4~km)
while the pointwise grid value deepens to $-3.1$, so the PARK
threshold was never formally crossed and the team held the core in
the core-seek TRACK mode throughout, a smoothing effect to fold into threshold
selection for the full tensor-line comparison, which remains future
work.

\begin{figure}[!t]
    \centering
    \includegraphics[width=0.48\textwidth]{figures/ocean_oecs_overlay_2km.png}
    % Generated by experiments/main_ocean_hfr_2km_oecs.py, under
    % trunk/Python_Simulations/Vector_Fields/VF_Robot/. Review copy in
    % ocean_data/det_jacobian_plots/oecs_trajectory_overlay_2km.png.
    \caption{$s_1$ tracker (core-seek) on the 2~km Santa Barbara Channel
    field over the 24-h forward FTLE background. White stars are TRAP
    cores extracted offline from the full radar grid at the final
    frame; the team converges from the start (green) to a core near the
    channel entrance and station-keeps for the rest of the 28-h record,
    ending 0.8~km from the nearest offline core.}
    \label{fig:ocean_oecs}
\end{figure}


% ---------------------------------------------------------------------------
\section{Discussion}
\label{sec:discussion}

\subsection{Comparison with Prior Work}
\label{sec:prior_comparison}

In prior work \cite{14}, a three-robot cluster fit a linear model to each
component of $\mathbf{v}$, assembled the estimated Jacobian
$\mathbf{J}$ and offset $\mathbf{h}$, and recovered the critical point
$\mathbf{p}^* = -\mathbf{J}^{-1}\mathbf{h}$ where $\mathbf{v}(\mathbf{p}^*) =
\mathbf{0}$, together with its type from the eigenvalues of $\mathbf{J}$.
A linear model, however, carries no information about how $\mathbf{J}$
itself varies in space. The structures addressed in this paper, boundaries
between rotation-dominated and strain-dominated flow and the separatrices
attached to saddle points, are level sets and valley lines of $D$, and
navigating them requires the first and second spatial derivatives of $D$:
a field-structure problem in place of a point-estimation problem, and the
reason the estimation burden rises from three robots to six.

\subsection{Implications for Distributed Eulerian-Structure Following}

The stability result of Section~\ref{sec:stability} is harder to
obtain than a codimension-one regular-level-set tracking result would
be, for a structural reason. A regular level set is determined by
first-order information alone: a nonvanishing gradient fixes the
transversal direction, and a controller built from it admits a global
certificate from a single quadratic Lyapunov function. The trench
lives exactly where that first-order information vanishes in the
normal direction, so controller and proof alike must synthesize the
missing direction from curvature; the eigenstructure dependence, mode
switching, and tube assumptions of Section~\ref{sec:stab_sep} are the
price of that synthesis, and the Hessian estimator of
Section~\ref{sec:estimation}, rather than the gradient estimate a
first-order feature would need, is the equipment. What the results add
is that a six-robot formation sampling only its own local velocities
recovers enough second-order structure to pay that price, with no map,
no forecast, and no communication beyond the single synchronous sample
each control cycle already requires.

\subsection{Park versus Traverse: One Threshold, Two Behaviors}
\label{sec:park_vs_traverse}

Theorem~\ref{thm:separatrix} and Remark~\ref{rem:capture_gap} identify
the terminal-capture test $\lambda_1\lambda_2 \geq 0$ as the reason the
deployed controller continues through every saddle in
Section~\ref{sec:results}: the biased $\hat{\mathbf{H}}_D$ makes that
test unreliable exactly where it would need to fire. An estimator-aware
alternative restores capture without touching the transverse Newton
snap, the flow step, or the slide step, using only the two quantities
Section~\ref{sec:sensitivity} shows are linear in the fitted
coefficients and free of the Hessian's structural bias, $\hat{D}$ and
$\hat{\mathbf{g}}$. The \emph{park step} adds a fourth mode, checked
before the FLOW test:
\begin{equation}
    \text{if } \hat{D} < D_{\text{capture}}:\quad
    \dot{\mathbf{p}}_c^{\,\text{des}} =
    -\begin{bmatrix}
        \mathrm{sat}(\hat{g}_x) \\ \mathrm{sat}(\hat{g}_y)
    \end{bmatrix},
    \label{eq:park_step}
\end{equation}
componentwise saturated gradient descent on $\hat{D}$, which carries
its own certificate: with exact estimates $D$ itself is the Lyapunov
function, $\dot{D} = -\sum_k g_k\,\mathrm{sat}(g_k) \leq 0$ with
equality only at critical points, so the team settles at the nearest
local minimum of $D$, the flow saddle at the end of the
current trench segment. $D_{\text{capture}}$ is the one-parameter
knob: $D_{\text{capture}} = -\infty$ disables the test and the
controller is trajectory-identical to the base selector (verified
bit-exact in \textit{separatrix\_logic\_park\_step}), while a finite
threshold inside the well (here $-0.5$, roughly half the double-gyre
well depth $\pi^4 A^2 \approx 0.974$) parks the team there.
Fig.~\ref{fig:park_vs_traverse} runs both settings from the same start,
noise-free: traverse exits the domain at $(1.00, -0.51)$ after 228
steps, riding the crossing trench; park reaches the saddle and holds
position, settling 0.0124~m from $\mathbf{p}^*_1$ with zero
tail-position standard deviation over its final 100 steps. This is the
practical resolution of Remark~\ref{rem:capture_gap}: continuation is
the honest behavior of the theorem's capture test under a real
estimator, and one threshold recovers whichever behavior a mission
calls for, terminal parking at a critical point or open-ended traversal
of the network.

\begin{figure}[!t]
    \centering
    \includegraphics[width=0.45\textwidth]{figures/park_vs_traverse.png}
    % Generated by experiments/park_vs_traverse_demo.py, under
    % trunk/Python_Simulations/Vector_Fields/VF_Robot/. Review copy in
    % experiments/outputs/separatrix_clean/.
    \caption{Park versus traverse from the same start $(0.10, -0.20)$,
    noise-free. Traverse (the default selector) rides the trench through
    the bottom saddle and exits the domain. Park
    ($D_{\text{capture}} = -0.5$) applies the park step
    (\ref{eq:park_step}) once $\hat{D}$ drops below threshold and holds
    the team 0.0124~m from the saddle.}
    \label{fig:park_vs_traverse}
\end{figure}

\subsection{What Objectivity Buys, and What It Costs}
\label{sec:objectivity_tradeoff}

The two controllers are not competitors so much as the two sides of
the tradeoff carried by (\ref{eq:splitting}). The $D$ tracker exploits
everything the instantaneous velocity field offers, including the flow
direction itself: its flow-assisted mode carries it through the crest,
and its trench network is globally connected, so one team can traverse
an entire separatrix system in a single run. The price, made concrete
by Fig.~\ref{fig:oecs_objectivity}, is the spin term: a rotating
observer, a drifting reference, or any platform whose navigation frame
turns relative to the Earth inflates $\omega^2/4$ as in
(\ref{eq:D_not_objective}) and tracks a curve no material particle
follows. The $s_1$ tracker keeps only the strain term, and what it
recovers is the guarantee that matters for material prediction: the
curve it parks on is where drifting objects actually collect, in any
frame, which is precisely the search-and-rescue use case TRAPs were
built for \cite{36}. What it gives up follows from the same refusal:
bounded segments that end at degeneracies, so no crest crossing and no
network traversal; gradient-rate transverse convergence, since
Remark~\ref{rem:s1hess} shows no quadratic fit can supply the
curvature a Newton snap needs (the measured zero-noise tracking error,
0.008 against the $D$ tracker's 0.002, is this cost made visible); and a
tangent oriented by continuity, since a direction field admits no
global sign. A mission that needs long-range structure traversal in a
known, fixed frame is better served by the $D$ tracker, and one that needs
frame-proof material prediction, or operates from a rotating or
drifting platform, by the $s_1$ tracker. Both consume the same fit, so
the selection is one function call, not a new estimator or formation.

\subsection{Estimator Diagnosis}
\label{sec:estimator_diagnosis}

Open-loop checks at the nominal radius confirm the structural picture
of Section~\ref{sec:s1_estimation}. The fitted shear strain at three
test points is at the $10^{-4}$ level (analytic value zero); the fitted
$\hat{\mathbf{H}}_{s_1}$ eigenvalues are $[-0.011, 0.000]$ and smaller,
negative semidefinite as Remark~\ref{rem:s1hess} requires and
negligible against the analytic transverse curvature of $6.9$--$9.7$;
and the transverse component of (\ref{eq:grad_s1}) recovers the
restoring slope $4.0$ per unit offset. Under measurement noise the
first-order advantage of the eigenframe is an order of magnitude: at
$\sigma_{uv} = 0.01$, the median direction error of $\mathbf{e}_2$ is
$4.5^\circ$, against $44^\circ$ for $\nabla\hat{s}_1$ and $41^\circ$
for the $\hat{\mathbf{H}}_D$ eigenvector that the $D$ tracker uses as its
tangent at the same point ($2\times 10^4$ draws per level). The
tracker's tangent therefore remains informative far beyond the noise
level at which both gradient channels approach chance, which is the
mechanism behind the noise-robustness margins of
Section~\ref{sec:results_oecs}.

\subsection{Failure Modes}

Three behaviors fall outside the guarantees of
Section~\ref{sec:stability} and are worth stating plainly. First, the
estimator is ill-conditioned whenever the six sampling points drift
toward a common conic (Proposition~\ref{prop:conic}):
$\det(\boldsymbol{\Phi})$ shrinks continuously as the formation
deforms toward one, until at the conic itself the fit is undefined. A
ring without a center robot is the standing example, and a controller
could in principle drive the formation toward a degenerate shape;
monitoring $\det(\boldsymbol{\Phi})$ and adding a shape-maintenance
term to the centroid command is the direct mitigation. The Monte Carlo
trials flag the weaker symptom, a root-mean-square inter-robot
distance error exceeding 0.30 (four times the formation radius), as a
formation collapse, and it does not occur in any reported trial.

Second, the separatrix controller can chatter across the FLOW-band
boundary (\ref{eq:flow_band}) when noise pushes $\hat{D}$ or the
dimensionless ratio back and forth across the threshold. Because every
mode shares the transverse dynamics (\ref{eq:n_dynamics}), chatter
does not destabilize the transverse channel; its cost is a lower
effective along-trench speed. $\varepsilon_{\text{dim}}$ trades this
against the crest crossing of Section~\ref{sec:dg_example}: too small
a threshold reintroduces the crest stall, too large widens the chatter
band. The value used throughout, $\varepsilon_{\text{dim}} = 0.025$,
was fixed once and never retuned.

Third, the $s_1$ tracker's eigenframe is noise-fragile exactly where
$\mathbf{S}$ approaches isotropy: the angular error of the frame
scales as the coefficient noise divided by the eigenvalue gap $2r$, a
property of the physics rather than of the estimator, since the tensor
lines themselves are singular there
(Section~\ref{sec:strain_field}). The selector's
$\varepsilon_{\text{degen}}$ guard hands these regions to the ACQUIRE
mode, which uses only $\nabla\hat{s}_1$; the residual risk is a noisy
tangent-sign flip when a trial passes close to a degeneracy, which the
continuity rule damps but cannot exclude, and which the generic ocean
field, whose isolated degeneracies move between frames, exercises
directly.


\subsection{Limitations}
\label{sec:limitations}

All results are from simulation; hardware validation on the Decabot
testbed is planned as a follow-on study. The double-gyre results use
the static field ($\varepsilon = 0$): in the time-varying case the
separatrix becomes a moving material curve, and although the ocean
experiments show both trackers operating on a genuinely time-varying
field without modification, no analytic guarantee is offered for that
case. Beyond the double gyre, generalization is demonstrated only on
the ocean dataset, though the estimation framework applies to any
smooth planar vector field. Finally, the stability guarantees, and the
reported success rates, are local. Theorem~\ref{thm:separatrix}
assumes initialization inside the tube $T_w$, and all Monte Carlo
trials start from a formation straddling the separatrix; the attract
mode converges to the critical set of $D$, which includes the
gyre-center maxima, so a team started deep inside a vortex core parks
at the core instead. In supplementary uniform-random-start trials over
the domain, roughly half (47.9\%, $n = 1000$) of starts reached the
trench at zero noise. We do not characterize the basin of attraction
of the separatrix; prior adaptive-navigation primitives \cite{27},
\cite{28} likewise report local convergence without basin estimates.


% ---------------------------------------------------------------------------
\section{Conclusion and Future Work}
\label{sec:conclusion}
\label{sec:future}

This paper presented one distributed estimator and two trackers built
on it. A six-robot formation, proved minimal, recovers the local
quadratic model of a planar flow from a single synchronous sample, and
with it both terms of the splitting identity $D = \omega^2/4 - s_1^2$,
the Jacobian determinant and the smaller strain eigenvalue, with
closed-form, heading-isotropic noise gains verified to $0.3\%$. The
$D$ tracker rides the connected trench network of $D$ on a
modified Newton step, with a Lyapunov certificate and a
network-traversal theorem whose terminal-capture case is one threshold
away. The $s_1$ tracker rides the bounded attracting segments of $s_1$;
its design follows from a structural result, that no quadratic velocity
fit can certify $s_1$-trench curvature, and its frame equivariance was
demonstrated in closed loop rather than asserted: the same material
curve in either observer frame (gap 0.020), where the $D$ tracker
follows a frame artifact (gap 1.219). Monte Carlo sweeps placed the
$D$ tracker's failure cliff exactly at the estimator's
gradient-reliability threshold (measured-to-predicted ratio 1.0), so
the estimator, not the control law, sets its noise ceiling; the $s_1$
tracker's first-order eigenframe held its 50\% level to ten times that
noise. On real Santa
Barbara Channel currents the $D$ tracker followed a coherent
ridge for 28 hours to a consistent landfall reproduced by 84 of 100
nearby starts, and the $s_1$ tracker converged to and held a transient
attracting profile within one radar grid cell of ground truth computed
offline from the full map. Every capability rests on the same premise:
a team that can sense second-order structure locally does not need the
map.

Four extensions are direct. A time-derivative correction term to the
navigation primitive, analyzed against the analytic $\varepsilon > 0$
double gyre where the exact structure motion is known, would extend
the stability results to time-varying fields with a quantified bound
rather than empirical evidence alone. Online formation reshaping,
adapting the pair lengths and SAS parameters, would protect the
conditioning of $\boldsymbol{\Phi}$ near degenerate configurations.
Hardware deployment on the Decabot testbed would close the loop on
physical robots, with the field emitted by an onboard computer and
sampled by each robot's own sensors. And extension to three
dimensions raises the estimation stakes: a full 3D quadratic model
carries ten coefficients per component, so ten robots, or a
reduced-order estimator with different assumptions.


\appendices

% ---------------------------------------------------------------------------
\section{Double-Gyre Analytic Forms}
\label{app:analytic}

This appendix derives the closed-form field quantities of the
double-gyre running example (Section~\ref{sec:dg_example}), used as
ground truth throughout the paper.

For the steady double gyre with amplitude $A$, $x_f = x + 1$,
$y_f = y + 0.5$, define $s_x = \sin(\pi x_f)$, $c_x = \cos(\pi x_f)$,
$s_y = \sin(\pi y_f)$, $c_y = \cos(\pi y_f)$. The field is
$u = -\pi A s_x c_y$, $v = \pi A c_x s_y$.

The Jacobian entries are
\begin{equation}
    \frac{\partial u}{\partial x} = -\pi^2 A c_x c_y, \quad
    \frac{\partial u}{\partial y} =  \pi^2 A s_x s_y,
\end{equation}
\begin{equation}
    \frac{\partial v}{\partial x} = -\pi^2 A s_x s_y, \quad
    \frac{\partial v}{\partial y} =  \pi^2 A c_x c_y,
\end{equation}
so the flow is divergence free and
\begin{equation}
    \begin{split}
    D = \det(\mathbf{J})
      &= -\pi^4 A^2 \bigl(c_x^2 c_y^2 - s_x^2 s_y^2\bigr) \\
      &= -\frac{\pi^4 A^2}{2}\bigl(\cos 2\pi x_f + \cos 2\pi y_f\bigr),
    \end{split}
    \label{eq:det_analytic}
\end{equation}
using the half-angle identities. Differentiating
(\ref{eq:det_analytic}),
\begin{equation}
    \begin{aligned}
    \nabla D &= \pi^5 A^2
    \begin{bmatrix} \sin 2\pi x_f \\ \sin 2\pi y_f \end{bmatrix},
    \\
    \mathbf{H}_D &= 2\pi^6 A^2
    \begin{bmatrix} \cos 2\pi x_f & 0 \\ 0 & \cos 2\pi y_f \end{bmatrix}.
    \end{aligned}
    \label{eq:gradhess_analytic}
\end{equation}
The Hessian is diagonal everywhere, so $D$ is additively separable in
$x_f$ and $y_f$, which is what makes the trench normal form of
Assumption~\ref{ass:trench} exact for this field. The vorticity is
$\omega = -2\pi^2 A s_x s_y$, which vanishes on the separatrix
$x_f = 1$. The shear strain is $\tfrac{1}{2}(\pi^2 A s_x s_y - \pi^2 A
s_x s_y) = 0$, so the rate-of-strain tensor is diagonal everywhere,
$\mathbf{S} = \mathrm{diag}(u_x, -u_x)$ with $u_x = -\pi^2 A c_x c_y$.
Its eigenvalues are $\pm|u_x|$, giving (\ref{eq:s1_analytic_main}),
with eigenvectors on the coordinate axes: the compression direction is
$\hat{\mathbf{x}}$ where $u_x < 0$ and $\hat{\mathbf{y}}$ where $u_x >
0$, the segment-identity alternation of Section~\ref{sec:dg_example}.
$\mathbf{S}$ is degenerate exactly where $c_x c_y = 0$, the lines
$x_f \in \tfrac{1}{2} + \mathbb{Z}$ and $y_f \in \tfrac{1}{2} +
\mathbb{Z}$.


% ---------------------------------------------------------------------------
\section{Pentagon Cluster Kinematics}
\label{app:kinematics}

This appendix defines the twelve-state cluster parameterization used
by the cluster space controller of Section~\ref{sec:arch}.

The pentagon-of-pairs cluster arranges six robots as three pairs; pair
$k$ ($k = 1, 2, 3$) consists of robots $2k-1$ and $2k$ with midpoint
$\mathbf{m}_k = (\mathbf{p}_{2k-1} + \mathbf{p}_{2k})/2$, half-length
$L_k/2$, and orientation $\psi_k$, so its robots sit at $\mathbf{m}_k
\pm (L_k/2)(\cos\psi_k, \sin\psi_k)^\top$. The three midpoints form a
triangle parameterized by the SAS variables of the three-robot case in
\cite{14}: $p_1 = \|\mathbf{m}_2 - \mathbf{m}_1\|$, $q_1 =
\|\mathbf{m}_3 - \mathbf{m}_2\|$, interior angle $\beta_1$ at
$\mathbf{m}_2$, centroid $\mathbf{p}_c = (\mathbf{m}_1 + \mathbf{m}_2 +
\mathbf{m}_3)/3$, and heading $\theta_c = \text{atan2}(m_{2,y} -
m_{1,y}, m_{2,x} - m_{1,x})$. The full 12-state cluster configuration
is $(x_c, y_c, \theta_c, p_1, \beta_1, q_1, L_1, \psi_1, L_2, \psi_2,
L_3, \psi_3)$: the six triangle states plus the length and orientation
of each of the three pairs, all free, matching the twelve degrees of
freedom of six planar robots.
% Code naming offset: pentagon_kinematics.py calls the pair variables
% L_2..L_4, theta_2..theta_4 for pairs A..C (no L_1/theta_1 exists in
% code); the math is identical.
The inverse kinematics reconstructs the midpoint triangle
from $(p_1, \beta_1, q_1)$ in a local frame, rotates and translates it
by $(\theta_c, x_c, y_c)$, and offsets each pair from its midpoint,
following Appendix B of \cite{14}; differentiating analytically gives
the $12 \times 12$ inverse kinematic Jacobian that maps desired cluster
state velocities to the twelve robot velocity commands. The explicit
entries are omitted; the implementation is in
\textit{pentagon\_kinematics.py}.


% ---------------------------------------------------------------------------
\section{Proofs}
\label{app:proofs}

This appendix collects the proofs of the estimation results of
Section~\ref{sec:estimation} and the stability results of
Section~\ref{sec:controllers}. Statements are restated where they were
summarized rather than given formally in the body.

\begin{proposition}[Minimality]
\label{prop:minimality}
Each robot contributes exactly one linear equation to each system in
(\ref{eq:solve}), and each system has six unknowns. Hence no algorithm,
linear or otherwise, can uniquely recover the quadratic model
(\ref{eq:quad_model_u})--(\ref{eq:quad_model_v}) from one synchronous
sample of fewer than six robots, and six robots suffice exactly when
$\boldsymbol{\Phi}$ is nonsingular.
\end{proposition}

\begin{IEEEproof}
With $N < 6$ robots, $\boldsymbol{\Phi} \in \mathbb{R}^{N \times 6}$ has
a nontrivial null space, so for any candidate coefficient vector
$\boldsymbol{\theta}_u$ the vector $\boldsymbol{\theta}_u +
\boldsymbol{\theta}_0$ with $\boldsymbol{\Phi}\boldsymbol{\theta}_0 =
\mathbf{0}$, $\boldsymbol{\theta}_0 \neq \mathbf{0}$, generates identical
measurements at every robot. The two quadratic fields are
indistinguishable from the data, so no estimator can select between
them. With $N = 6$, (\ref{eq:solve}) has a unique solution if and only if
$\det(\boldsymbol{\Phi}) \neq 0$.
\end{IEEEproof}

\begin{proposition}[Conic criterion]
\label{prop:conic}
$\boldsymbol{\Phi}$ is singular if and only if the six sampling points
lie on a common conic, that is, on the zero set of some nonzero quadratic
polynomial $q(\tilde{x}, \tilde{y})$ (including degenerate conics such as
line pairs).
\end{proposition}

\begin{IEEEproof}
$\boldsymbol{\Phi}$ is singular iff its columns are linearly dependent,
iff there is a nonzero coefficient vector $\mathbf{c}$ with
$\boldsymbol{\phi}(\tilde{\mathbf{p}}_i)^\top \mathbf{c} = 0$ for all
$i$. The function $q = \boldsymbol{\phi}^\top \mathbf{c}$ is precisely a
nonzero polynomial of degree at most two vanishing at all six points.
\end{IEEEproof}

\begin{corollary}[Pentagon plus center]
\label{cor:pentagon}
Place five robots at the vertices of a regular pentagon of radius $\rho$
and the sixth at its center. Then $\boldsymbol{\Phi}$ is nonsingular.
\end{corollary}

\begin{IEEEproof}
No three vertices of a regular pentagon are collinear, and five points in
general position determine a unique conic; for the pentagon vertices that
conic is the circumscribed circle. The center does not lie on the
circle, so no conic passes through all six points, and
Proposition~\ref{prop:conic} applies.
\end{IEEEproof}

\begin{IEEEproof}[Proof of Lemma~\ref{lem:gains}]
The value coefficient is exact: every basis function except the constant
vanishes at the origin, so interpolation forces $\hat{u}_0$ to equal the
center robot's reading, giving unit gain. The remaining rows follow from
direct inversion of $\boldsymbol{\Phi}$ using the fifth roots of unity;
the row norms reduce to the stated constants, and repeating the
computation with an arbitrary ring phase leaves them unchanged. The
error covariance of the estimated gradient is in fact isotropic,
$\tfrac{4}{10}\sigma_{uv}^2\,\rho^{-2}\mathbf{I}$, for every heading.
\end{IEEEproof}

\begin{IEEEproof}[Proof of Theorem~\ref{thm:attract}]
Since $\nabla V_A = \mathbf{H}_D \mathbf{g}$,
$\dot{V}_A = \mathbf{g}^\top \mathbf{H}_D \boldsymbol{\delta}
= -\sum_k v_{\max} \tanh(z_k/v_{\max})\,
\mathbf{g}^\top \mathbf{H}_D \mathbf{w}_k
= -\sum_k v_{\max} \tanh(z_k/v_{\max})\,\lambda_k
(\mathbf{w}_k^\top\mathbf{g})$, and $\lambda_k(\mathbf{w}_k^\top
\mathbf{g}) = \lambda_k^2 z_k$, giving (\ref{eq:VA_dot}); each summand
is nonnegative because $z\tanh z \geq 0$. Equality requires $z_1 = z_2 =
0$, that is $\mathbf{g} = \mathbf{0}$. Convergence to the critical set
follows from the invariance principle. Near an isolated critical point
$|z_k| \ll v_{\max}$, the saturation is inactive at first order,
$\boldsymbol{\delta} \approx -\mathbf{H}_D^{-1}\mathbf{g} \approx
-(\mathbf{p}_c - \mathbf{p}^\star_D)$, a unit-rate linear contraction.
\end{IEEEproof}

\begin{IEEEproof}[Proof of Lemma~\ref{lem:transverse}]
By (\ref{eq:normal_form}) the Hessian is diagonal in
$(\hat{\mathbf{e}}_\parallel, \hat{\mathbf{e}}_\perp)$ with transverse
eigenvalue $\lambda_\perp > 0$ and eigenvector
$\hat{\mathbf{e}}_\perp$, and the transverse gradient component is
$\hat{\mathbf{e}}_\perp^\top \mathbf{g} = \lambda_\perp n$. The
$\hat{\mathbf{e}}_\perp$ component of the attract step
(\ref{eq:attract_step}), of the slide step (\ref{eq:slide_step}) (the
transverse eigenvalue is positive, so the signed and magnitude
denominators agree), and of the flow step (\ref{eq:flow_step}) are all
built from the same argument $-\lambda_\perp n / \lambda_\perp = -n$.
Saturation gives (\ref{eq:n_dynamics}). (For the double gyre itself the
normal form holds to second order rather than exactly: the exact
transverse gradient and curvature give the Newton argument
$-\tan(2\pi n)/(2\pi)$ in place of $-n$, which has the sign of $n$ and
magnitude at least $|n|$ for $|n| < 1/4$, so the resulting contraction
is at least as fast as (\ref{eq:n_dynamics}) and every conclusion that
follows is unchanged.) The right side is negative for $n > 0$ and
positive for $n < 0$, so $V_\perp = \tfrac{1}{2}n^2$ is a common
Lyapunov function for the switched system; for $|n| \leq v_{\max}$,
$\tanh(n/v_{\max}) \geq n \tanh(1)/v_{\max}$ yields $\dot{V}_\perp \leq
-2\tanh(1) V_\perp$.
\end{IEEEproof}

\begin{IEEEproof}[Proof of Lemma~\ref{lem:progress}]
On $\Gamma$ the tangential gradient component is $\varphi'(s)$ and the
along-trench eigenvector is $\hat{\mathbf{e}}_\parallel$. In the FLOW band
the tangential command is $v_{\max}\tanh((\mathbf{f}^\top
\mathbf{w}_1)/v_{\max})$ applied along $\mathbf{w}_1 = \pm
\hat{\mathbf{e}}_\parallel$; the product form makes it a nonnegative
multiple of the flow's tangential component, and by
Assumption~\ref{ass:flow} the crest region is covered by this mode, with
magnitude at least the first argument of (\ref{eq:progress_rate}). Off
the band, the slide step is selected precisely when the flow agrees with the
descent direction $-\mathrm{sign}(\varphi')\hat{\mathbf{e}}_\parallel$,
and its tangential command moves along descent with magnitude
$v_{\max}\tanh(|\varphi'|/(|\lambda_1| v_{\max}))$; the attract step is selected
when the flow opposes descent, and its tangential command,
$-\varphi'/\lambda_1$ with $\lambda_1 < 0$, moves along ascent, which by
the selector's own condition is again the flow direction, with the same
magnitude. In the attract fallback ($\mathbf{H}_D \succeq 0$ near the
wells) the tangential command descends toward the well, which is the
flow direction near a stable-manifold endpoint. In all cases the sign
agrees with the flow and the magnitude is bounded below by
(\ref{eq:progress_rate}) via Assumption~\ref{ass:flow}.
\end{IEEEproof}

\begin{IEEEproof}[Proof of Theorem~\ref{thm:separatrix}]
(i) is Lemma~\ref{lem:transverse}. (ii): by Lemma~\ref{lem:progress},
$s(t)$ is monotone in the flow direction with rate at least $m$ outside
the $r_0$ balls, and mode switches cannot reverse it, so the ball is
reached in time at most $L_\Gamma/m$. Capture (iii): inside the ball,
$\mathbf{H}_D \succ 0$, the selector is in its attract fallback, and
Theorem~\ref{thm:attract} gives local exponential convergence to the
well, which coincides with the flow saddle $\mathbf{p}^*$. Continuation:
if the fallback test fails to trigger despite $\mathbf{H}_D \succ 0$
holding for the true field, the selector instead falls through to
Lemma~\ref{lem:progress}'s flow or descent case on whichever of the two
crossing segments the flow direction and sign convention of
$\hat{\mathbf{g}}$ select, and (i)-(ii) apply unchanged with $\Gamma$
relabeled.
\end{IEEEproof}

\begin{IEEEproof}[Proof of Proposition~\ref{prop:iss}]
$\tfrac{d}{dt}\tfrac{1}{2}n^2 = n(-v_{\max}
\tanh(n/v_{\max}) + d) < 0$ whenever $v_{\max}\tanh(|n|/v_{\max}) >
\bar{d}$, that is outside the stated interval, which is therefore
attractive and invariant.
\end{IEEEproof}

\begin{IEEEproof}[Proof of Proposition~\ref{prop:equivariance}]
$\mathbf{S}' = \mathbf{Q}\mathbf{S}\mathbf{Q}^\top$, so $s_1' = s_1$,
$r' = r$, $\mathbf{e}_k' = \mathbf{Q}\mathbf{e}_k$, and $\nabla' s_1' =
\mathbf{Q}\nabla s_1$ (the gradient of an invariant scalar co-rotates).
The branch conditions are scalar thresholds on the invariants $r$,
$\hat{s}_1$, and the banded flag, so each branch fires identically in
both frames. The TRACK channels apply saturations of the invariant
inner products
($\nabla' s_1'^\top \mathbf{e}_k' = \nabla s_1^\top \mathbf{e}_k$)
along the co-rotating $\mathbf{e}_k$, and TRIM outputs zero, so both
output the rotated command $\mathbf{Q}\boldsymbol{\delta}$. The ACQUIRE
and PARK bodies saturate the two components of $\nabla\hat{s}_1$ in the
observer's coordinate axes; componentwise saturation commutes with
$\mathbf{Q}$ only where it is inactive, so these branches are exactly
covariant only in their linear regime, while under saturation each
command component still opposes the corresponding gradient component,
giving $\nabla\hat{s}_1^\top \boldsymbol{\delta} < 0$ and descent of
the invariant $\hat{s}_1$ in either frame. The one non-covariant input, the flow
$\mathbf{f}$ used to seed the ride tangent at the first tracking step,
enters only through the sign of one inner product at one instant;
thereafter orientation is propagated by continuity, which is itself
covariant.
\end{IEEEproof}


% ---------------------------------------------------------------------------
\begin{thebibliography}{99}

\bibitem{1} J. Cortés and M. Egerstedt, "Coordinated control of multi-robot
systems: A survey," \textit{SICE Journal of Control, Measurement, and
System Integration}, vol. 10, no. 6, pp. 495--503, 2017.

\bibitem{2} J. Das et al., "Coordinated sampling of dynamic oceanographic
features with underwater vehicles and drifters," \textit{The International
Journal of Robotics Research}, vol. 31, no. 5, pp. 626--646, 2012.

\bibitem{3} S. McCammon et al., "Ocean front detection and tracking using a
team of heterogeneous marine vehicles," \textit{Journal of Field Robotics},
vol. 38, no. 6, pp. 854--881, 2021.

\bibitem{4} S. R. Ramp et al., "Preparing to predict: the second autonomous
ocean sampling network (AOSN-II) experiment in the Monterey Bay,"
\textit{Deep Sea Research Part II: Topical Studies in Oceanography}, vol. 56,
no. 3-5, pp. 68--86, 2009.

\bibitem{5} C. Dong et al., "Oceanic mesoscale eddies," \textit{Ocean-Land-
Atmosphere Research}, vol. 4, p. 0081, 2025.

\bibitem{6} G. Haller and G. Yuan, "Lagrangian coherent structures and mixing
in two-dimensional turbulence," \textit{Physica D: Nonlinear Phenomena},
vol. 147, no. 3-4, pp. 352--370, 2000.

\bibitem{7} T. Peacock and G. Haller, "Lagrangian coherent structures: The
hidden skeleton of fluid flows," \textit{Physics Today}, vol. 66, no. 2,
pp. 41--47, 2013.

\bibitem{8} S. C. Shadden, F. Lekien, and J. E. Marsden, "Definition and
properties of Lagrangian coherent structures from finite-time Lyapunov
exponents in two-dimensional aperiodic flows," \textit{Physica D}, vol. 212,
no. 3-4, pp. 271--304, 2005.

\bibitem{9} M. Michini, M. A. Hsieh, E. Forgoston, and I. B. Schwartz,
"Robotic tracking of coherent structures in flows," \textit{IEEE Transactions
on Robotics}, vol. 30, no. 3, pp. 593--603, 2014.

\bibitem{10} M. Michini, H. Rastgoftar, M. A. Hsieh, and S. Jayasuriya,
"Distributed formation control for collaborative tracking of manifolds in
flows," in \textit{Proc. American Control Conference}, pp. 3874--3880, 2014.

\bibitem{11} M. Michini, M. A. Hsieh, E. Forgoston, and I. B. Schwartz,
"Experimental validation of robotic manifold tracking in gyre-like flows,"
in \textit{Proc. IEEE/RSJ International Conference on Intelligent Robots
and Systems (IROS)}, pp. 2306--2311, 2014.

\bibitem{12} G. Knizhnik, P. Li, X. Yu, and M. A. Hsieh, "Flow-based control
of marine robots in gyre-like environments," in \textit{Proc. 2022 Int. Conf.
on Robotics and Automation (ICRA)}, pp. 3047--3053, 2022.

\bibitem{13} D. Kularatne and M. A. Hsieh, "Tracking attracting manifolds in
flows," \textit{Autonomous Robots}, vol. 41, no. 8, pp. 1575--1588, 2017.

\bibitem{14} C. Waight and C. A. Kitts, "Adaptive navigation of multirobot
systems to critical points in 2D vector fields," \textit{IEEE/ASME Transactions
on Mechatronics}, to appear.

\bibitem{15} L. Briñón-Arranz, A. Renzaglia, and L. Schenato, "Multirobot
symmetric formations for gradient and Hessian estimation with application to
source seeking," \textit{IEEE Transactions on Robotics}, vol. 35, no. 3,
pp. 782--789, 2019.

\bibitem{16} G. Haller, "Lagrangian coherent structures," \textit{Annual
Review of Fluid Mechanics}, vol. 47, pp. 137--162, 2015.

\bibitem{17} M. Serra and G. Haller, "Objective Eulerian coherent
structures," \textit{Chaos: An Interdisciplinary Journal of Nonlinear
Science}, vol. 26, no. 5, p. 053110, 2016.

\bibitem{18} P. J. Nolan, M. Serra, and S. D. Ross, "Finite-time Lyapunov
exponents in the instantaneous limit and material transport,"
\textit{Nonlinear Dynamics}, vol. 100, no. 4, pp. 3825--3852, 2020.

\bibitem{19} T. D. Harms, S. T. M. Dawson, and B. J. McKeon, "Lagrangian
gradient regression for the detection of coherent structures from sparse
trajectory data," \textit{Royal Society Open Science}, vol. 11, no. 9,
p. 240586, 2024.

\bibitem{20} G. Haller, "Dynamic rotation and stretch tensors from a
dynamic polar decomposition," \textit{Journal of the Mechanics and Physics
of Solids}, vol. 86, pp. 70--93, 2016.

\bibitem{21} A. Okubo, "Horizontal dispersion of floatable particles in the
vicinity of velocity singularities such as convergences," \textit{Deep-Sea
Research}, vol. 17, no. 3, pp. 445--454, 1970.

\bibitem{22} J. Weiss, "The dynamics of enstrophy transfer in two-dimensional
hydrodynamics," \textit{Physica D: Nonlinear Phenomena}, vol. 48, no. 2-3,
pp. 273--294, 1991.

\bibitem{23} J. C. R. Hunt, A. A. Wray, and P. Moin, "Eddies, streams, and
convergence zones in turbulent flows," in \textit{Proc. Summer Program,
Center for Turbulence Research}, Rep. CTR-S88, pp. 193--208, 1988.

\bibitem{24} B.-L. Hua and P. Klein, "An exact criterion for the stirring
properties of nearly two-dimensional turbulence," \textit{Physica D}, vol. 113,
no. 1, pp. 98--110, 1998.

\bibitem{25} R. Cooper et al., "Initial study of multirobot adaptive
navigation for exploring environmental vector fields," in \textit{Proc.
IEEE Aerospace Conference}, 2019.

\bibitem{26} P. Ögren, E. Fiorelli, and N. E. Leonard, "Cooperative control
of mobile sensor networks: Adaptive gradient climbing in a distributed
environment," \textit{IEEE Transactions on Automatic Control}, vol. 49, no. 8,
pp. 1292--1302, 2004.

\bibitem{27} R. T. McDonald, C. A. Kitts, and M. A. Neumann, "Navigation of
scalar fronts with multirobot clusters in simulation and experiment,"
\textit{IEEE Systems Journal}, vol. 14, no. 3, pp. 3755--3766, 2020.

\bibitem{28} C. A. Kitts, R. T. McDonald, and M. A. Neumann, "Adaptive
navigation control primitives for multirobot clusters: Extrema finding,
contour following, ridge/trench following, and saddle point station keeping,"
\textit{IEEE Access}, vol. 6, pp. 17625--17642, 2018.

\bibitem{29} R. T. McDonald, C. A. Kitts, and M. A. Neumann, "Experimental
implementation and verification of scalar field ridge, trench, and saddle
point maneuvers using multirobot adaptive navigation," \textit{IEEE Access},
vol. 7, pp. 62950--62961, 2019.

\bibitem{30} C. A. Kitts and I. Mas, "Cluster space specification and control
of mobile multirobot systems," \textit{IEEE/ASME Transactions on Mechatronics},
vol. 14, no. 2, pp. 207--218, 2009.

\bibitem{31} T. Adamek, C. A. Kitts, and I. Mas, "Gradient-based cluster space
navigation for autonomous surface vessels," \textit{IEEE/ASME Transactions on
Mechatronics}, vol. 20, no. 2, pp. 506--518, 2014.

\bibitem{32} S. Tomer et al., "A low-cost indoor testbed for multirobot
adaptive navigation research," in \textit{Proc. IEEE Aerospace Conference},
pp. 1--12, 2018.

\bibitem{33} C. Waight and C. Kitts, "A functional indoor testbed for
multirobot adaptive navigation in vector field environments," in
\textit{Proc. ASME Int. Design Engineering Technical Conf. and Computers
and Information in Engineering Conf.}, vol. 89251, 2025.

\bibitem{34} Y. Sun, G. Ma, M. Liu, C. Li, and J. Liang, "Distributed
finite-time coordinated control for multi-robot systems with disturbances
and directed topologies," \textit{Transactions of the Institute of
Measurement and Control}, vol. 39, no. 12, pp. 2907--2919, 2017.

\bibitem{35} Y. N. Dauphin, R. Pascanu, C. Gulcehre, K. Cho, S. Ganguli,
and Y. Bengio, "Identifying and attacking the saddle point problem in
high-dimensional non-convex optimization," in \textit{Advances in Neural
Information Processing Systems}, vol. 27, pp. 2933--2941, 2014.

\bibitem{36} M. Serra, P. Sathe, I. Rypina, A. Kirincich, S. D. Ross,
P. Lermusiaux, A. Allen, T. Peacock, and G. Haller, "Search and rescue
at sea aided by hidden flow structures," \textit{Nature Communications},
vol. 11, no. 1, p. 2525, 2020.

\end{thebibliography}


\begin{IEEEbiography}[{
\includegraphics[width=1in,height=1.25in,clip,keepaspectratio]{figures/cwaight.jpg}}]{Christopher Waight}
received his BSc Degree in Physics from Santa Clara University. He later earned his Masters in Mechanical Engineering there, and is now a PhD Candidate in Mechanical Engineering. He currently works as a Research Assistant in SCU's Robotics System Lab. He also has experience working as a hardware engineer and as a Technical Program Manager in R\&D.
\end{IEEEbiography}

\begin{IEEEbiography}[{
\includegraphics[width=1in,height=1.25in,clip,keepaspectratio]{figures/ckitts.jpg}}]{Christopher A. Kitts}
(Senior Member, IEEE) earned degrees from Princeton University (B.S.E.), University of Colorado (M.P.A.), and Stanford University (M.S. and Ph.D.). He served as a Satellite Constellation Mission Controller with the U.S. Air Force and as a Computer Scientist at NASA Ames Research Center. He is currently Professor of Mechanical Engineering, Director of the Robotic Systems Laboratory, and Associate Dean at Santa Clara University. Dr. Kitts is a Fellow of ASME.
\end{IEEEbiography}

\end{document}
